
%%%%%%%%%%  scan idx 8  |  folio 1  %%%%%%%%%%


\underline{INTERSECTIONS SUR LES SURFACES REGULIERES}


\underline{par P. Deligne}

Dans la première partie de cet exposé, après quelques préliminaires de théorie des intersections (contenus dans $[5] \cup [7]$), on donne la démonstration d'après Artin d'un cas particulier d'un théorème de Raynaud [9] (1.13), et une démonstration a priori d'un théorème que Néron [8] vérifiait cas par cas (1.15). De (1.13) résultent assez aisément les théorèmes de [9] utilisés dans [2] et par là dans IX.

Dans la seconde partie, je poursuis sur ma lancée et raconte comment les ``italiens'' [1], [3] calculaient les nombres d'intersections, sur les surfaces. Les résultats principaux 2.9 et 2.11 se trouvent par ailleurs dans Northcott [12] et [13]. Pour terminer, je donne en 2.27, 2.30, 2.34 une relation conjecturale entre la dimension du schéma formel versel des modules d'une singularité d'une courbe réduite, l'invariant $\delta$ de la singularité et un invariant mesurant sa symétrie.

Cet exposé ne sera pas utilisé dans la suite du séminaire.


1. \underline{Nombres d'intersection sur les surfaces arithmétiques}

1.0. On suppose donné dans ce \S{} un schéma régulier $X$ de dimension 2, un diviseur réduit $Y$ de $X$, et un sous-corps $k$ de $H^0(Y,\underline{O}_Y)$. On suppose que $Y$ est une courbe (EGA II 7.4.1) propre sur $k$. Voici trois exemples de cette situation

\underline{Exemple} 1.1. Soit $(S,\eta,s)$ un trait et $X$ un $S$-schéma $f : X \longrightarrow S$, régulier, propre et plat sur $S$, à fibres purement de dimension un. On prend $Y = (f^{-1}(s))_{red}$ et $k = k(s)$. On se restreindra à ce cas à partir de 1.10.


%%%%%%%%%%  scan idx 9  |  folio 2  %%%%%%%%%%

\underline{Exemple} 1.2. Soient $X_o$ le spectre d'un anneau local normal de dimension 2, $s$ le point fermé de $X_o$ et $f : X \longrightarrow X_o$ un $X_o$-schéma régulier et connexe, propre sur $X_o$, distinct de $X_o$ mais coïncidant avec $X_o$ au-dessus de $X_o \setminus \{s\}$. On prend $Y = (f^{-1}(s))_{red}$ et $k = k(s)$.

\underline{Exemple} 1.3. On prend pour $X$ une surface régulière complète sur un corps $k$ et pour $Y$ une courbe réduite tracée sur $X$.

Rappelons le lemme suivant (SGA 6 IV ex. 1.9).

\underline{Lemme} 1.4. \underline{Soit} $Z_o$ \underline{un sous-schéma fermé d'un schéma noethérien} $Z$. \underline{Alors, l'inclusion de} $Z_o$ \underline{dans} $Z$ \underline{induit un isomorphisme du groupe de Grothendieck de la catégorie des faisceauxcchérents de} $\underline{O}_{Z_o}$ \underline{-modules dans le groupe de Grothendieck de la catégorie des faisceaux cohérents sur} $Z$ \underline{à support dans} $Z_o$ .

Ce lemme permet de définir la caractéristique d'Euler-Poincaré $\chi(\mathcal{F})$ d'un faisceau cohérent $\mathcal{F}$ sur $X$ à support dans $Y$, de façon à ce que,

(i) Si $\mathcal{F}$ est un faisceau cohérent de $\underline{O}_Y$-modules, alors

\begin{equation}\tag{1.4.1}
\chi(\mathcal{F}) \;=\; \dim_k H^0(\mathcal{F}) - \dim_k H^1(\mathcal{F})
\end{equation}
(ii) $\chi$ est additif relativement aux suites exactes courtes.

Si les faisceaux de cohomologie d'un complexe $K$ de $\underline{O}_X$-modules sont des faisceaux cohérents à support dans $Y$, on pose

\begin{equation}\tag{1.4.2}
\chi(K) \;=\; \Sigma \, (-1)^i \, \chi(\underline{H}^i(K)) \qquad .
\end{equation}
Si $\mathcal{F}$ est localisé en un point $y$ de $Y$, on a

\begin{equation}\tag{1.4.3}
\chi(\mathcal{F}) \;=\; \mathrm{lg}_{\underline{O}_y}(\mathcal{F}_y) \, . \, [k(y) : k] \qquad .
\end{equation}

%%%%%%%%%%  scan idx 10  |  folio 3  %%%%%%%%%%

Sous les hypothèses (1.1), (1.2) ou (1.3), on a

\begin{equation}\tag{1.4.4}
\chi(\mathcal{F}) \;=\; \mathrm{lg}\; f_{*}\mathcal{F} \;-\; \mathrm{lg}\; R^{1}f_{*}\mathcal{F} \quad .
\end{equation}
Soient $D$ un diviseur $\geq 0$ concentré sur $Y$, et $\mathcal{L}$ un faisceau inversible sur $D$. Il y a alors une et une seule façon de définir le degré $\deg_D(\mathcal{L})$ de sorte que

\begin{equation}\tag{1.4.5}
\left\{
\begin{array}{l}
\deg_D(\mathcal{L}' \otimes \mathcal{L}'') \;=\; \deg_D(\mathcal{L}') + \deg_D(\mathcal{L}'') \\[3ex]
\text{pour } \mathcal{L} \subset \underline{O} \text{ , on a } \deg_D(\mathcal{L}) \;=\; - \chi(\underline{O}/\mathcal{L}) \quad .
\end{array}
\right.
\end{equation}
On a alors la formule de Riemann-Roch sur $D$

\begin{equation}\tag{1.4.6}
\chi(\mathcal{L}) \;=\; \chi(\underline{O}_D) \;+\; \deg_D(\mathcal{L}) \quad .
\end{equation}
Soit $D = \sum n_i \, D_i$ la décomposition de $D$ en diviseurs irréductibles réduits ; on a

\begin{equation}\tag{1.4.7}
\deg_D(\mathcal{L}) \;=\; \sum n_i \, \deg_{D_i}(\mathcal{L}) \quad .
\end{equation}
\underline{Définition} 1.5. \underline{Soient} $D$ \underline{et} $E$ \underline{deux diviseurs} $\geq 0$ \underline{sur} $X$, \underline{et supposons} $D$ \underline{concentré sur} $Y$. \underline{On définit le nombre d'intersection de} $D$ \underline{et} $E$ \underline{(relativement au corps} $k$ (1.0)) \underline{par} :

\begin{equation*}
(D,E) \;=\; \chi(\underline{O}_D \overset{\mathbb{L}}{\otimes} \underline{O}_E) \;=\; \chi(\underline{O}_D \otimes \underline{O}_E) \;-\; \chi(\mathrm{Tor}^{1}(\underline{O}_D,\underline{O}_E)) \quad .
\end{equation*}
Sous les hypothèses faites, on a en effet $\mathrm{Tor}^{2}(\underline{O}_D,\underline{O}_E) = 0$ .

\underline{Proposition} 1.6. \underline{Sous les hypothèses} 1.5, \underline{on a} :

(i) \quad \underline{Le nombre d'intersection} $(D,E)$ \underline{est additif en les diviseurs} $D$ \underline{et} $E$ ; \underline{ceci permet}, \underline{par linéarité}, \underline{de définir le nombre d'intersection de deux diviseurs non nécessairement positifs.}


%%%%%%%%%%  scan idx 11  |  folio 4  %%%%%%%%%%

(ii) \quad $(D,E) \;=\; \deg_D(\underline{O}(E))$.

(iii) \quad \underline{Si} $D$ \underline{et} $E$ \underline{ne se coupent qu'en un nombre fini de points}, \underline{alors} $(D,E) = \chi(\underline{O}_D \otimes \underline{O}_E) \geq 0$, \underline{l'égalité ne pouvant avoir lieu que si} $D$ \underline{et} $E$ \underline{sont disjoints.}

(iv) \quad \underline{Si} $D$ \underline{et} $E$ \underline{sont concentrés sur} $Y$, \underline{alors} $(D,E) = (E,D)$.

(v) \quad \underline{Si} $N$ \underline{est le faisceau conormal à} $D$, \underline{alors}

\begin{equation*}
(D)^{2} \;=\; (D,D) \;=\; - \deg_D(N) \quad .
\end{equation*}
(vi) \quad \underline{Si} $E$, \underline{non nécessairement positif}, \underline{est le diviseur d'une fonction méromorphe sur} $X$, \underline{alors} $(D,E) = 0$.

\underline{Preuve}. Soit la suite exacte

\begin{equation*}
0 \longrightarrow \underline{O}(-E) \longrightarrow \underline{O}_X \longrightarrow \underline{O}_E \longrightarrow 0 \quad ,
\end{equation*}
elle fournit, par (1.4.6),

\begin{equation*}
\begin{array}{rl}
(D,E) \;\overset{\mathrm{dfn}}{=\!=}\; & \chi(\underline{O}_D \overset{\mathbb{L}}{\otimes} \underline{O}_E) \;=\; \chi(\underline{O}_D \overset{\mathbb{L}}{\otimes} \underline{O}_X) \;-\; \chi(\underline{O}_D \overset{\mathbb{L}}{\otimes} \underline{O}(-E)) \\[3ex]
= & \chi(\underline{O}_D) \;-\; \chi(\underline{O}(-E)\vert D) \;=\; - \deg_D (\underline{O}(-E)) \;=\; \deg_D (\underline{O}(E)) \quad .
\end{array}
\end{equation*}
Ceci prouve (ii), (vi) et la biadditivité en vertu de (1.4.5) et (1.4.7). Sous les hypothèses (iii), on a $\mathrm{Tor}^{1}(\underline{O}_D,\underline{O}_E) = 0$, d'où la conclusion. Sous les hypothèses (v), on a

\begin{equation*}
\underline{O}_D \otimes \underline{O}_D \;=\; \underline{O}_D \quad \text{et} \quad \mathrm{Tor}^{1}(\underline{O}_D,\underline{O}_D) \;\simeq\; N \quad (\text{SGA 6 VII 2.5}) \quad ,
\end{equation*}
d'où

\begin{equation*}
(D,D) \;=\; \chi(\underline{O}_D) \;-\; \chi(N) \;=\; - \deg_D(N) \quad .
\end{equation*}
Enfin, (iv) résulte de la symétrie du produit tensoriel.

On notera que (ii), (iv) et (vi) restent vrais sans supposer les diviseurs positifs.


%%%%%%%%%%  scan idx 12  |  folio 5  %%%%%%%%%%

\underline{Proposition} 1.7. \underline{Soient} $(Y_i)_{1 \leq i \leq k}$ \underline{les composantes irréductibles de} $Y$, $f$ \underline{une fonction méromorphe sur} $X$, \underline{et}

\begin{equation*}
\mathrm{div}(f) \;=\; \sum n_i\, Y_i \;+\; Z
\end{equation*}
\underline{son diviseur}, \underline{avec} $|Z| \cap Y$ \underline{fini}. \underline{Alors}, \underline{posant} $D_i = n_i\, Y_i$, \underline{on a}, \underline{pour} $(a_i) \in \mathbb{Q}^k$ :

\begin{equation*}
(\sum a_i\, D_i)^2 \;=\; -\sum_{i<j} (a_i - a_j)^2\, (D_i, D_j) \;-\; \sum a_i^2\, (D_i, Z) \qquad .
\end{equation*}
\underline{Preuve} (d'après [7]). En vertu de 1.6 (v) (vi), on a

\begin{equation*}
\begin{aligned}
(\sum_i a_i\, D_i)^2 \;&=\; \sum_i a_i\, (D_i,\; \sum_j a_j\, D_j) \;=\; \sum_i a_i\, (D_i,\; \sum_j a_j\, D_j - a_i\ \mathrm{div}\ f) \\[2pt]
&=\; \sum a_i\, (D_i,\; \sum_{j \neq i} (a_j - a_i) D_j \;-\; a_i Z) \\[2pt]
&=\; \sum_{i \neq j} a_i\, (a_j - a_i)(D_i, D_j) \;-\; \sum_i a_i^2\, (D_i, Z) \\[2pt]
&=\; -\sum_{i<j} (a_i - a_j)^2\, (D_i\ D_j) \;-\; \sum_i a_i^2\, (D_i, Z)
\end{aligned}
\end{equation*}
\underline{Corollaire} 1.8. \underline{Sous les hypothèses de l'exemple} 1.1, \underline{prenons pour} $f$ \underline{une uniformisante} $\pi$ \underline{de} $S$. \underline{Ici}, $Z = 0$ \underline{et}

\begin{equation*}
(\sum k_i\, Y_i)^2 \;=\; -\sum_{i<j} \frac{1}{n_i n_j}\; (k_i\, n_j - k_j\, n_i)^2\, (Y_i,\, Y_j) \qquad ;
\end{equation*}
\underline{donc la forme d'intersection est négative}, \underline{et si} $Y$ \underline{est connexe}, \underline{son noyau est réduit aux multiples de} $\mathrm{div}(\pi)$.

\underline{Corollaire} 1.9 (Mumford [7]). \underline{Sous les hypothèses de l'exemple} 1.2, \underline{prenons pour} $f$ \underline{un élément non nul de l'idéal maximal de l'anneau local} $\underline{O}_a$ . \underline{Ici}, $Z > 0$, \underline{de sorte que la forme d'intersection est définie négative}.


%%%%%%%%%%  scan idx 13  |  folio 6  %%%%%%%%%%

1.10. \underline{On se restreint}, \underline{pour la suite}, \underline{au cas considéré dans l'exemple} (1.1).

Lichtenbaum [5] a prouvé que, dans ce cas, $X$ est automatiquement projectif sur $S$, de sorte qu'il n'est pas fatiguant de définir le complexe $Rf^{!}\underline{O}_S$ ([4] III 8.7 p. 190). Puisque $f$ est de dimension relative 1 et que $X$ et $S$ sont réguliers(Gorenstein suffirait) il existe ([4] V 8.3 et V 9.1) un faisceau inversible $\omega_{X/S}$ sur $S$, unique à isomorphisme unique près, tel que

\begin{equation}\tag{1.10.1}
Rf^{!}\underline{O}_S \;\simeq\; \omega_{X/S}[1] \quad .
\end{equation}
On désigne par $K$ un quelconque diviseur (effectif ou non) tel que $\omega \simeq \underline{O}(K)$.

Soient $D$ un diviseur $\geq 0$ concentré sur $Y$ et $N$ son faisceau conormal :

\begin{center}
\begin{tikzcd}[column sep=large, row sep=large]
D \arrow[rr, "i"] \arrow[dr, swap, "g"] & {} & X \arrow[dl, "f"] \\
 {} & S &
{} \end{tikzcd}
\end{center}
% STRUCTURE: Hand-drawn triangle, exactly 3 nodes and 3 arrows. NODES: (1) D, upper left; (2) X, upper
%   right; (3) S, bottom centre (apex pointing down). ARROWS: (a) D --> X, horizontal, pointing RIGHT,
%   single plain arrowhead at X, label "i" set ABOVE the shaft; (b) D --> S, diagonal descending to the
%   RIGHT, single plain arrowhead at S, label "g" set to the LEFT of the shaft (outside the triangle);
%   (c) X --> S, diagonal descending to the LEFT, single plain arrowhead at S, label "f" set to the
%   RIGHT of the shaft (outside the triangle). No arrow is hooked/mono, epi/twoheaded, dashed, or marked
%   iso; none carries a tilde. The thickened ink at the D-end of arrow (b) is the start-blob of the
%   hand-drawn stroke, verified at high zoom to be NOT a second arrowhead. A lone full stop is typed on
%   the baseline of S, well to the RIGHT of the diagram (sentence punctuation); it is not a node and is
%   omitted from the tikzcd.
On a

\begin{equation}\tag{1.10.2}
Rg^{!}\underline{O}_S = Ri^{!}Rf^{!}\underline{O}_S = Ri^{!}\omega_{X/S}[1] = \omega_{X/S}\big|\ D \otimes N^{\vee}[0]
\end{equation}
En vertu du théorème de dualité ([4] V 11.1 p.210),

\begin{equation}\tag{1.10.3}
Rg_{*}\, Rg^{!}\underline{O}_S \;\simeq\; \mathrm{RHom}(Rg_{*}\, \underline{O}_D,\underline{O}_S) \quad .
\end{equation}
De plus, pour tout $\underline{O}_S$-module $M$ concentré au point fermé, on a

\begin{equation}\tag{1.10.4}
\chi(\mathrm{RHom}(M,\underline{O}_S)) \;=\; -\,\chi(M)
\end{equation}
de sorte qu'ici, utilisant (1.4.4) et (1.10.3), on obtient

\begin{equation*}
\chi(Rg^{!}\underline{O}_S) \;=\; -\,\chi(\underline{O}_D) \quad .
\end{equation*}
Puisque (1.4.6)

\begin{equation*}
\chi(Rg^{!}\underline{O}_S) = \chi(\underline{O}_D) + \deg_D(Rg^{!}\underline{O}_S) \qquad , 
\end{equation*}

%%%%%%%%%%  scan idx 14  |  folio 7  %%%%%%%%%%

 on en conclut par $(1.10.2)$ que

\begin{equation*}
- 2\,\chi(\underline{O}_D) \;=\; \deg_D\,(Rg^{!}\underline{O}_D) \;=\; \deg_D\,(\omega|D \otimes N^{\vee}) \;=\; \deg_D\,(\omega) - \deg_D(N) \qquad ,
\end{equation*}
ce qui, par 1.6 (ii) et(v) se réécrit

\underline{Proposition} 1.11. \underline{Sous les hypothèses} (1.1), \underline{quelque soit le diviseur} $D \geq 0$ \underline{concentré sur} $Y$, \underline{on a}

\begin{equation*}
\chi(\underline{O}_D) \;=\; - \frac{1}{2}\,(D,D+K) \qquad .
\end{equation*}
1.12. On désignera par $d$ le pgcd des multiplicités géométriques des composantes irréductibles de $X_s = f^{-1}(s)$, i.e. le pgcd des longueurs des anneaux locaux des points maximaux de la fibre spéciale géométrique de $f$. Si $X_{\eta}$ a un point rationnel, i.e. si $f : X \longrightarrow S$ a une section $x$, alors $d = 1$ ; plus précisément $f$ est lisse en $x(s)$. En effet, $x$ est une immersion fermée de schémas réguliers, donc est une immersion régulière et on applique EGA IV 17.12.1.

Voici la démonstration, d'après M. Artin, d'un cas particulier d'un théorème de M. Raynaud [9], qui suppose seulement $d$ premier à la caractéristique $p$ de $k$.

\underline{Théorème} 1.13. \underline{Avec les notations précédentes}, \underline{si} $Y$ \underline{est géométriquement connexe et si} $d = 1$, \underline{alors} $X$ \underline{est cohomologiquement plat sur} $S$ \underline{en dimension} 0 (EGA III 7.8.1), \underline{plus précisément} (EGA III 7.6.9 (ii)) \underline{on a}

\begin{equation*}
k \xrightarrow{\;\sim\;} H^0(X_s,\, \underline{O}_{X_s}).
\end{equation*}
On se ramène aisément au cas où $S$ est strictement local. Comme $Y$ est géométriquement connexe, $H^0(Y,\underline{O}_Y)$ est une extension radicielle de $k$, et puisque $(d,p)=1$, une au moins des composantes irréductibles de $Y$ est
%%%%%%%%%%  scan idx 15  |  folio 8  %%%%%%%%%%
géométriquement réduite, et on en déduit facilement que $H^0(Y,\underline{O}_Y) = k$.

Soit $D$ un diviseur sur $X$ tel que

\begin{equation*}
Y \leq D < X_s \qquad .
\end{equation*}
Puisque $d = 1$, $D$ n'est pas multiple rationnel de $X_s$, donc par 1.8

\begin{equation*}
(D,X_s-D) \;=\; - \,(D,D) \;>\; 0 \qquad .
\end{equation*}
Le diviseur $Y$ a donc une composante irréductible $E$ telle que

\begin{equation*}
0 < E \leq X_s-D \quad \text{et} \quad (D,E) > 0 \qquad .
\end{equation*}
Considérons le diagramme

\begin{center}
\begin{tikzcd}[column sep=large, row sep=large]
0 \arrow[r] & \underline{O}(-D-E) \arrow[r] \arrow[d, "\varphi"] & \underline{O} \arrow[r] \arrow[d, equal] & \underline{O}_{D+E} \arrow[r] \arrow[d, "\psi"] & 0 \\
0 \arrow[r] & \underline{O}(-D) \arrow[r] & \underline{O} \arrow[r] & \underline{O}_D \arrow[r] & 0
\end{tikzcd}
\end{center}
% STRUCTURE: Two rows of five nodes each, 10 nodes total, arranged in a 2x5 grid. TOP ROW (left to
%   right): node T1 = 0 ; node T2 = \underline{O}(-D-E) ; node T3 = \underline{O} ; node T4 =
%   \underline{O}_{D+E} ; node T5 = 0. BOTTOM ROW (left to right): node B1 = 0 ; node B2 =
%   \underline{O}(-D) ; node B3 = \underline{O} ; node B4 = \underline{O}_D ; node B5 = 0. HORIZONTAL
%   ARROWS, top row (4, all plain single-headed \longrightarrow pointing right, all unlabelled): T1 ->
%   T2 ; T2 -> T3 ; T3 -> T4 ; T4 -> T5. HORIZONTAL ARROWS, bottom row (4, all plain single-headed
%   \longrightarrow pointing right, all unlabelled): B1 -> B2 ; B2 -> B3 ; B3 -> B4 ; B4 -> B5. VERTICAL
%   ARROWS (3, all pointing downward from the top row to the bottom row): (i) T2 -> B2, plain
%   single-headed arrow pointing down, labelled \varphi, the label printed to the RIGHT of the arrow
%   shaft; (ii) T3 -> B3, NOT an arrow: two parallel vertical strokes with no arrowhead at either end,
%   i.e. an equality/identity bar between \underline{O} and \underline{O}; unlabelled; (iii) T4 -> B4,
%   plain single-headed arrow pointing down, labelled \psi, the label printed to the RIGHT of the arrow
%   shaft. No arrows are hooked, doubled-headed, dashed or marked iso. Columns 1 and 5 (the zeros) carry
%   no vertical arrows. OUTSIDE THE DIAGRAM: a comma is printed at the far right of the page, on the
%   baseline of the bottom row, after node B5; it is the terminal punctuation of the displayed diagram
%   and cannot be placed inside the tikzcd body.
le diagramme du serpent fournit un isomorphisme entre $\mathrm{Ker}(\psi)$ et $\mathrm{Coker}(\varphi)$. Ce conoyau n'est autre que la restriction de $\underline{O}(-D)$ à $E$, d'où une suite exacte

\begin{equation*}
0 \longrightarrow \underline{O}(-D)|E \longrightarrow \underline{O}_{D+E} \longrightarrow \underline{O}_D \longrightarrow 0 \qquad .
\end{equation*}
Par construction $\deg_E(\underline{O}(-D) = - (D,E) < 0$, donc $H^0(\underline{O}(-D)|E) = 0$, et la flèche de restriction

\begin{equation*}
0 \longrightarrow H^0(\underline{O}_{D+E}) \longrightarrow H^0(\underline{O}_D)
\end{equation*}
est injective. Ceci permet de construire par récurrence une suite de diviseurs

\begin{equation*}
Y = D_0 \leq D_1 \leq \ldots \leq D_n = X_s
\end{equation*}
telle que $H^0(\underline{O}_{D_i})$ s'injecte dans $H^0(\underline{O}_{D_{i-1}})$, donc $H^0(\underline{O}_{X_s}) = H^0(\underline{O}_Y) = k$, ce qui achève la démonstration de 1.13.


%%%%%%%%%%  scan idx 16  |  folio 9  %%%%%%%%%%

Le fait suivant est prouvé dans [5] ou [10].

\underline{Théorème} 1.14. (Castelnuovo, Enriques,...). \underline{Sous les hypothèses} (1.1), \underline{soit} $D$ \underline{un diviseur réduit irréductible contenu dans} $Y$. \underline{Supposons que} $H^1(D,\underline{O}_D) = 0$ \underline{et que,} \underline{posant} $k' = H^0(D,\underline{O}_D)$, \underline{on ait}

\begin{equation*}
(D,D) \; = \; - [k' : k] \qquad .
\end{equation*}
\underline{Alors,} $D$ \underline{est une droite projective sur} $k'$ \underline{et il existe un} $S$-\underline{schéma régulier} $X_1$ , \underline{un point fermé} $x_1$ \underline{de} $X_1$ \underline{de corps résiduel} $k'$ \underline{et un isomorphisme entre} $X$ \underline{et le} $S$-\underline{schéma déduit de} $X_1$ \underline{par éclatement de} $x_1$, \underline{cet isomorphisme identifiant} $D$ \underline{à la fibre de} $x_1$ \underline{dans} $X_1$.

On appelle un diviseur $D$ possédant ces propriétés une \underline{courbe exceptionnelle de première espèce}.

Je ne résiste pas au plaisir de démontrer la proposition suivante de Néron [8] prop 1' p. 90 :

\underline{Proposition} 1.15. \underline{Supposons que} $X_\eta$ \underline{soit une courbe elliptique} \underline{(en particulier ait une section)} \underline{et que} $X$ \underline{ne contienne aucune courbe exceptionnelle de première espèce.} \underline{Alors,} \underline{avec la notation} (1.10.1), \underline{on a}

\begin{equation*}
f^*f_* \; \omega_{X/S} \qquad \stackrel{\sim}{\longrightarrow} \qquad \omega_{X/S}
\end{equation*}
\underline{universellement}.

Comme $d = 1$ (1.12), il suffit, d'après (1.13), de prouver que $\omega_{X/S}$ est isomorphe à $\underline{O}_X$ ; puisque $\omega_{X/S}$ est trivial sur la courbe elliptique $X_\eta$ , on sait a priori que

\begin{equation*}
\omega_{X/S} \qquad \simeq \qquad \underline{O}(\Sigma \, n_i \, Y_i) \qquad , 
\end{equation*}

%%%%%%%%%%  scan idx 17  |  folio 10  %%%%%%%%%%

 où les $Y_i$ sont les composantes irréductibles de $Y$, et puisque $\mathrm{div}(\pi) \sim 0$, on peut supposer les $n_i$ positifs. Si $\Sigma \, n_i \, Y_i$ est proportionnel au diviseur $\mathrm{div}(\pi)$, il est multiple entier de $\mathrm{div}(\pi)$ et on a gagné. Sinon, on a

\begin{equation*}
(\Sigma \, n_i \, Y_i)^2 < 0 \qquad ,
\end{equation*}
il existe $i$ tel que

\begin{equation*}
(Y_i, \; \Sigma \, n_j \, Y_j) < 0 \qquad ,
\end{equation*}
et on applique le lemme suivant :

\underline{Lemme} 1.16. \underline{Sous les hypothèses} 1.1, \underline{supposons que} $Y$ \underline{soit connexe et que} $Y$ \underline{soit réductible,} \underline{ou que} $\chi(\underline{O}_Y) \leq 0$. \underline{Alors,} \underline{pour qu'une composante irréductible} $Y_i$ \underline{de} $Y$ \underline{soit une courbe exceptionnelle de première espèce,} \underline{il faut et il suffit que} $(Y_i,K) < 0$ .

Soit $k' = H^0(Y_i,\underline{O}_{Y_i})$. La formule de Riemann-Roch 1.11 fournit

\begin{equation*}
\chi(\underline{O}_{Y_i}) = [k':k] \; (1-\dim_{k'} H^1(Y_i,\underline{O}_{Y_i})) = - \frac{1}{2} \; ((Y_i,Y_i) + (Y_i,K)) \qquad .
\end{equation*}
Si $Y_i$ est exceptionnelle de première espèce, on a donc $(Y_i,K) = -[k':k] < 0$ . Réciproquement, si $(Y_i,K) < 0$, alors $\chi(\underline{O}_{Y_i}) > 0$, ce qui implique que $\chi(\underline{O}_{Y_i}) = [k':k]$ et que $Y$ soit réductible, de sorte que par 1.8, $(Y_i,Y_i) < 0$. Les nombres négatifs $(Y_i,Y_i)$ et $(Y_i,K)$ sont divisibles par $[k':k]$, de sorte qu'on doit avoir

\begin{equation*}
(Y_i,Y_i) = (Y_i,K) = -[k':k] \qquad .
\end{equation*}
\underline{Exercice} (a). Sous les hypothèses (1.1), on suppose que $X_\eta$ est un espace principal homogène sous une courbe elliptique et que $X$ ne contienne aucune courbe exceptionnelle de première espèce. Montrer que sur $X$, la formule de Riemann-Roch (1.11) s'écrit

\begin{equation*}
\chi(\underline{O}_D) = - \frac{1}{2} \, (D,D)
\end{equation*}

%%%%%%%%%%  scan idx 18  |  folio 11  %%%%%%%%%%

(b) Supposons en outre le corps résiduel $k(s)$ algébriquement clos. Montrer que le diagramme des composantes irréductibles de $X_s$ est multiple de l'un des diagrammes considérés par Néron [8] p. 124-125, 4è colonne (cas a à c 8). Ce diagramme détermine les genres des composantes irréductibles réduites, leur multiplicité dans $X_s$, la nature et la disposition des singularités de $(X_s)_{\mathrm{red}}$ .

On traitera d'abord le cas où $X_s$ est irréductible ; ce cas exclus, on prouvera que toute composante irréductible réduite est une droite projective de self-intersection $-$ 2. Si $X_s = \Sigma\ n_i\ D_i$ est la décomposition de $X_s$ en diviseurs irréductibles réduits, on réécrira cette condition

\begin{equation}\tag{*}
n_i \;=\; \frac{1}{2} \sum_{j \neq i} n_i \ (D_i \ D_j) \qquad\qquad .
\end{equation}
D'autre part, $X_s$ est connexe, de sorte que :

(**) Quels que soient $i$ et $j$, il existe une chaîne $i = i_1, i_2 \ldots i_n = j$ avec $(D_{i_k} \ D_{i_{k+1}}) \neq 0$ .

On cherchera alors toutes les solutions en nombre entiers de $(*) + (**)$. On montrera d'abord que, sauf dans les cas b2 et c2 de Néron, les nombres d'intersection $(D_i D_j)$ sont égaux à zéro ou un.


2. \underline{Intersections italiennes}

Soient $S'$ et $S$ deux surfaces projectives lisses sur un corps $k$ et $f : S' \longrightarrow S$ un morphisme birationnel. Si $D$ et $E$ sont deux diviseurs sur $S$, on a

\begin{equation*}
(D,E) = (f^* \, D, \ f^* \, E) \qquad\qquad , \qquad 
\end{equation*}

%%%%%%%%%%  scan idx 19  |  folio 12  %%%%%%%%%%

 comme on le voit, par exemple, en faisant bouger $D$ et $E$.

Il s'impose donc de donner une expression ``birationnelle'' pour le nombre d'intersection. On se propose d'expliquer comment procédaient les italiens pour ce faire. Ma principale source de renseignements est [1] où se trouve une longue bibliographie. Voir aussi Manin [6], et Northcott[12],[13].

(2.0) Soient $S$ le spectre d'un anneau local régulier de dimension 2, $s$ le point fermé de $S$ et $k$ un sous-corps de $k(s)$ dont $k(s)$ soit une extension finie. Les nombres d'intersections et les caractéristiques d'Euler-Poincaré seront calculés en terme de $k$.

Dans ce \S, un $S$-schéma $p : S' \longrightarrow S$ sera dit être un \underline{transformé birationnel de} $S$ s'il est régulier, connexe, et si $p$ est propre et induit un isomorphisme de $S' \setminus p^{-1}(s)$ sur $S \setminus \{s\}$. Tout ce qui sera dit de $S$ s'appliquera encore aux localisés de ses transformés birationnels en un point fermé.

(2.1) Soit $f : (S'',p'') \longrightarrow (S',p')$ un morphisme de $S$-schémas entre transformés birationnels de $S$. On sait alors [5] qu'il existe une et une seule factorisation de $f$ :

\begin{equation*}
f = f_n \ \ldots \ f_1 \qquad , \qquad f_i : S_{i+1} \longrightarrow S_i
\end{equation*}
telle que $S_{i+1}$ se déduise de $S_i$ par éclatement d'un ensemble fini $X_i$ de points fermés, et que $f_i(X_{i+1}) \subset X_i$ . L'ensemble $X_i$ est l'ensemble des points de $S_i$ en lesquels la projection de $S''$ sur $S_i$ n'est pas un isomorphisme.

En particulier, $f$ induit un isomorphisme au-dessus du complémentaire d'une partie finie de $S'$, qui contient les points génériques des composantes irréductibles de $p'^{-1}(s)$. Dès lors, $f^{-1}$ définit une injection
%%%%%%%%%%  scan idx 20  |  folio 13  %%%%%%%%%%
de l'ensemble $\mathrm{Irr}({p'}^{-1}(s))$ des composantes irréductibles de ${p'}^{-1}(s)$ dans $\mathrm{Irr}({p''}^{-1}(s))$.

Les transformés birationnels de $S$ n'ont pas de $S$-automorphisme . On désignera ici par $I$ l'ensemble des classes de $S$-isomorphie des transformés birationnels de $S$, ordonné par la relation de domination. L'ensemble $I$ est filtrant à droite, et de plus petit élément $S$.

\underline{Définition} 2.2. \underline{L'ensemble} $S^*$ \underline{des} points de $S$ infiniment voisins de $s$ \underline{est} \underline{l'ensemble limite inductive}

\begin{equation*}
S^* = \varinjlim_{I} \ \mathrm{Irr}\,(p^{-1}(s)) \qquad .
\end{equation*}
Si un élément $t \in S^*$ est représenté sur un transformé birationnel $S'$ de $S$ par une courbe complète réduite irréductible $T$, le corps $H^o(T,\underline{O}_T)$ ne dépend que de $t$, et non du choix de $S'$ ; on l'appelle le \underline{corps résiduel} $k(t)$ de $t$.

Soit $S_1$ l'éclaté de $S$ le long de $s$. La droite projective $s^*$ image réciproque de $\{s\}$ dans $S_1$ s'appelle la \underline{curvetta} de $s$ et définit un élément $s^*$ de $S^*$ de corps résiduel $k(s)$.

Soient $S'$ un transformé birationnel de $S$, $t$ un point fermé de $S'$ et $S'_t$ le spectre de l'anneau local de $S'$ en $t$. On définit de façon évidente une injection de ${S'}^*_t$ dans $S^*$ ; en particulier, $t$ définit un élément $t^*$ de $S^*$, de même corps résiduel. On désigne par ${S'}^*$ la réunion dans $S^*$ des ${S'}^*_t$ pour $t$ fermé dans $S'$.

La proposition suivante résulte de la description 2.1 des transformés birationnels de $S$ :


%%%%%%%%%%  scan idx 21  |  folio 14  %%%%%%%%%%

\underline{Proposition} 2.3. \underline{Quel que soit} $t \in S^*$, \underline{il existe une et une seule suite de transformés birationnels de} $S$ $(S_i)_{0 \leqslant i \leqslant n+1}$ \underline{et de points fermés} $t_i \in S_i$ $(0 \leqslant i \leqslant n)$ \underline{telle que}

(1)\quad $S_o = S$, \underline{et} $S_{i+1}$ \underline{est l'éclaté de} $S_i$ \underline{le long de} $\{t_i\}$

(ii)\quad $t_{i+1} \in t_i^*$\quad $(0 \leqslant i < n)$ \underline{et} $t = t_n^*$ .

\underline{De plus}, \underline{quel que soit le transformé birationnel} $(S',p')$ \underline{de} $S$, \underline{pour que} $t$ \underline{appartienne à l'image de} $\mathrm{Irr}({p'}^{-1}(s))$ \underline{dans} $S^*$, \underline{il faut et il suffit que} $S'$ \underline{domine} $S_{n+1}$.

On appelle $n$ l'ordre de $t$. On désignera ici par $S(t)$ (resp. par $S_+(t)$) le transformé birationnel $S_n$ (resp. $S_{n+1}$) de $S$. Par abus de notation, on désignera encore par $t$ le point fermé $t_n$ de $S_n$ et par $t^*$ son image réciproque $t_n^*$ dans $S_{n+1}$.

On munit $S^*$ de l'ordre suivant :

\begin{equation*}
t_1 \prec t_2 \quad \Longleftrightarrow \quad S_+(t_2) \ \mbox{domine} \ S_+(t_1) \qquad .
\end{equation*}
Si $t_1 \prec t_2$, alors $k(t_2)$ est une extension de $k(t_1)$. L'ensemble ordonné $S^*$ a pour plus petit élément $s^*$, et toute partie majorée de $S^*$ est finie et totalement ordonnée ($S^*$ est un ``arbre de souche $s^*$''). De plus, l'ordre $n$ de $t \in S^*$ est la distance, dans l'arbre $S^*$, de $s^*$ à $t$ ($n = \#[s^*,t]$).

La proposition suivante explicite la description 2.1 des transformés birationnels de $S$.

\underline{Proposition} 2.4. \underline{La fonction qui à un transformé birationnel} $p : S' \to S$ \underline{de} $S$ \underline{associe la partie} $\mathrm{Irr}(p^{-1}(s))$ \underline{de} $S^*$ \underline{induit une bijection entre l'ensemble} $I$ (2.1) \underline{et l'ensemble des parties finies de} $S^*$ \underline{qui}, \underline{avec tout élément},
%%%%%%%%%%  scan idx 22  |  folio 15  %%%%%%%%%%
\underline{contiennent ses prédécesseurs}.

Soient $D$ est un diviseur $\geq 0$ sur $S$ et $t \in S*$. On dit que $t$ \underline{appartient} à $D$ si le point fermé $t$ de $S(t)$ (2.3) appartient au transformé pur (= adhérence schématique de $D - \{s\}$) de $D$ sur $S(t)$.

\underline{Définition} 2.5. \underline{Soient} $t$, $t' \in S*$. \underline{On dit que} $t'$ \underline{est} proche \underline{de} $t$ \underline{si} $t' \in S_+(t)*$ (2.2) \underline{et appartient à la curvetta} $t*$ \underline{de} $t$ (2.3).

Si $t'$ est un successeur immédiat de $t$, alors $t'$ est proche de $t$ ; si dans une suite $(t_i)_{0 \leq i \leq n}$ chaque $t_{i+1}$ est successeur immédiat de $t_i$, et si $t_n$ est proche de $t_o$, alors $t_i$ est proche de $t_o$ pour $0 < i \leq n$ et $t_n$ n'est pas proche de $t_i$ pour $0 < i \leq n-1$.

(2.6)  Soient $A$ un anneau local noethérien de dimension $n$, $T$ son spectre, $m$ son idéal maximal, $k = A/m$ son corps résiduel et $P = \sum_{0 \leq i < n} a_i X^i$ son polynôme de Hilbert-Samuel :

\begin{equation*}
P(n) \;=\; \dim_k(m^n/m^{n+1}) \qquad \text{pour } n \text{ grand.}
\end{equation*}
Rappelons que la multiplicité $\mu_s$ dans $T$ du point fermé $s$ vérifie par définition

\begin{equation*}
a_{n-1} \;=\; \mu_s/(n-1)\, ! \qquad .
\end{equation*}
Lorsque $n = 1$, on a donc simplement

\begin{equation*}
\mu_s \;=\; \dim_k m^n/m^{n+1} \qquad \text{pour } n \text{ grand.}
\end{equation*}
En termes géométriques, $\mu_s$ peut se définir comme le degré, dans l'espace projectif sur $k$ $P_k(m/m^2)$, de la fibre spéciale de l'éclaté de $T$ le long du sous-schéma $\{s\}$.


%%%%%%%%%%  scan idx 23  |  folio 16  %%%%%%%%%%

Rappelons que si $Y$ et $Z$ sont deux sous-schémas fermés d'un schéma $X$, alors, le transformé pur de $Y$ dans l'éclaté $\tilde{X}$ de $X$ le long de $Z$, adhérence schématique de $Y \setminus Z$ dans $\tilde{X}$, n'est autre que l'éclaté de $Y$ le long de $Y \cap Z$.

Supposons que $T$ soit un sous-schéma fermé d'un schéma régulier $U$ . Soient $p : U_1 \longrightarrow U$ l'éclaté de $U$ le long de $\{s\}$, $s*$ l'image réciproque de $s$ dans $U_1$ et $T'$ le transformé pur de $T$ dans $U_1$. Ce qui précède montre que la multiplicité de $T$ en $s$ est le degré dans l'espace projectif $s*$ du sous-schéma $s* \cap T'$ .

Supposons en outre que $T$ soit un diviseur de $U$, soit $T*$ son transformé total dans $U_1$ (le diviseur image réciproque) et définissons $\nu$ par la formule $T* = \nu\, s* + T'$. Si $X$ est une droite de $s*$, et si on calcule les caractéristiques d'Euler-Poincaré comme sommes alternées de dimension sur $k$, on sait que

\begin{equation}\tag{2.6.1}
(X,s*) = \chi(\underline{O}_X \overset{\mathbb{L}}{\otimes} \underline{O}_{s*}) \;=\; \deg_X \underline{O}(s*) = -1
\end{equation}
et on a vérifié que

\begin{equation}\tag{2.6.2}
(X,T') = \chi(\underline{O}_X \overset{\mathbb{L}}{\otimes} \underline{O}_{T'}) \;=\; \deg_{s*} (s* \cap T') = \mu_s \qquad .
\end{equation}
La formule de projection

\begin{equation*}
Rp_*(\underline{O}_X \overset{\mathbb{L}}{\otimes} \underline{O}_{T*}) \;=\; Rp_*(\underline{O}_X \overset{\mathbb{L}}{\otimes} Lp^* \underline{O}_T) = Rp_* \underline{O}_X \overset{\mathbb{L}}{\otimes} \underline{O}_T
\end{equation*}
et son corollaire

\begin{equation*}
(X,T*) \;=\; \nu(X,s*) + (X,T') = 0
\end{equation*}
fournit, par (2.6.1), (2.6.2), l'identité

\begin{equation}\tag{2.6.3}
\mu_s \;=\; \nu \qquad .
\end{equation}

%%%%%%%%%%  scan idx 24  |  folio 17  %%%%%%%%%%

\underline{Proposition} 2.7. \underline{Soient} $T$ \underline{un diviseur} $\geq 0$ \underline{dans un schéma régulier} $U$, $s$ \underline{un point fermé de} $U$, $U_1$ \underline{l'éclaté de} $U$ \underline{le long de} $\{s\}$, $s^*$ \underline{l'image réciproque de} $s$ \underline{dans} $U_1$, $T'$ \underline{le transformé pur de} $T$ \underline{sur} $U_1$ \underline{et} $T^*$ \underline{son transformé total}.

\underline{Il y a identité entre}

(i) \underline{La multiplicité}, \underline{au sens de Samuel}, \underline{du point} $s$ \underline{de} $T$.

(ii) \underline{Le degré}, \underline{dans l'espace projectif} $s^*$, \underline{de} $s^* \cap T'$.

(iii) \underline{La multiplicité de} $s^*$ \underline{dans le diviseur} $T^*$.

2.8. Reprenons les hypothèses de 2.0. Soit $D$ un diviseur $\geq 0$ sur $S$ et $t \in S^*$. La \underline{multiplicité} $\mu_t(D)$ \underline{de} $D$ \underline{en} $t$ est, par définition, la multiplicité en $t$ du transformé pur de $D$ sur $S(t)$ (2.3). La fonction de $t$ $\mu_t(D)$ a pour support l'ensemble des points infiniment voisins de $s$ qui appartiennent à $D$ (2.4). Lorsque $D$ est régulier, les nombres $\mu_t(D).[k(t) : k(s)]$ sont égaux à zéro ou un : en effet, les transformés purs $D'$ de $D$ sont isomorphes à $D$, et on applique la définition (2.6) des multiplicités.

Les nombres d'intersections se calculent en termes de multiplicités :

\underline{Théorème} 2.9. \underline{Si les diviseurs} $D$ \underline{et} $E$ \underline{ne se coupent qu'en} $s$, \underline{on a}

\begin{equation*}
(D,E) \;=\; \sum_{t \in S^*} \mu_t(D)\mu_t(E)[k(t):k]
\end{equation*}
\underline{la somme étant étendue aux points infiniment voisins de} $s$ \underline{communs à} $D$ \underline{et} $E$.

\underline{Preuve}. Soit $p : S' \longrightarrow S$ un transformé birationnel de $S$. On désigne par $D^*$ et $E^*$ (resp. par $D'$ et $E'$) les transformés totaux (resp. purs) de $D$ et $E$.


%%%%%%%%%%  scan idx 25  |  folio 18  %%%%%%%%%%

\underline{Lemme} 2.9.1. (i) \underline{Le transformé total} $D^*$ \underline{est caractérisé par} :

(a) $D^* - D'$ \underline{est concentré sur} $f^{-1}(s)$

(b) \underline{si} $X$ \underline{est concentré sur} $f^{-1}(s)$, \underline{alors} $(D^*,X) = 0$

(ii) $(D,E) = (D^*,E^*) = (D^*,E') = (D',E^*)$.

L'assertion (i)(a) est triviale, et (i)(b) résulte de 1.6 (vi). Que (a) et (b) déterminent $D^*$ résulte de 1.9. Par construction, le conoyau de la flèche de $\underline{O}_E$ dans $f_* \underline{O}_E$ est fini, d'où

\begin{equation*}
(D^*,E') = \chi(f^*\underline{O}_D \otimes \underline{O}_{E'}) = \chi(\underline{O}_D \otimes f_* \underline{O}_{E'}) = \chi(\underline{O}_D \otimes \underline{O}_E) = (D,E)
\end{equation*}
ce qui, par (i), entraîne (ii).

Prenons pour $S'$ l'éclaté de $S$ en $s$. D'après (2.9.1), (2.6.3) et 2.7, on a

\begin{equation}\tag{2.9.2}
\begin{aligned} (D,E) &= (D^*,E') = (D' + \mu_s(S)\, s^*,E')\\ &= (D',E') + \mu_s(D)(s^*,E') = (D',E') + \mu_s(D)\mu_s(E)[k(s):k] \;. \end{aligned}
\end{equation}
Le théorème 2.9. se prouve maintenant par récurrence sur $(D,E)$, par une application itérée de (2.9.2).

(2.10) Soient $A$ un anneau local noethérien réduit de dimension un, $s$ le point fermé de $\mathrm{Spec}(A)$, $\tilde{A}$ le normalisé de $A$ et $k$ un corps dont $k(s)$ soit extension finie. Supposons que $\tilde{A}$ soit fini sur $A$ (i.e. que le complété $A^{\wedge}$ soit réduit). On appellera \underline{invariant} $\delta$ \underline{de la singularité de} $\mathrm{Spec}(A)$ \underline{en} $s$ la longueur, calculée sur $k$, du $A$-module $\tilde{A}/A$.


%%%%%%%%%%  scan idx 26  |  folio 19  %%%%%%%%%%

\underline{Théorème} 2.11. \underline{Soit} $D$ \underline{une courbe réduite tracée sur} $S$. \underline{Alors} :

(i) \underline{Quel que soit} $t \in S^*$, \underline{on a}, \underline{avec les notations} 2.8

\begin{equation*}
\mu_t(D) \;=\; \sum_{t' \text{ proche de } t} \mu_{t'}(D)[k(t'):k(t)] \qquad .
\end{equation*}
(ii) \underline{Supposons le normalisé} $\widetilde{D}$ \underline{de} $D$ \underline{fini sur} $D$, \underline{et soit} $\delta = \chi(\underline{O}_{\widetilde{D}}/\underline{O}_D)$ \underline{l'invariant de la singularité de} $D$ \underline{en} $s$. \underline{Pour tous les} $t \in S^*$, \underline{à l'exception d'un nombre fini}, $\mu_t(D)$ \underline{est égal à zéro ou un}, \underline{et on a}

\begin{equation*}
\delta \;=\; \sum_{t \in S^*} \frac{1}{2}\,\mu_t(D)\,(\mu_t(D)-1)\,[k(t):k] \qquad .
\end{equation*}
Soient $p : S_1 \longrightarrow S$ l'éclaté de $S$ le long de $\{s\}$, $D^*$ le transformé total de $D$ et $D'$ son transformé pur. En vertu de (2.7), on a

\begin{equation}\tag{2.11.1}
D^* = D' + \mu_s(D)\, s^* \qquad \text{et}
\end{equation}
\begin{equation}\tag{2.11.2}
\mu_s(D) = (s^*,D')\ .\ [k(s):k]^{-1} \qquad .
\end{equation}
La curvetta $s^*$ est non singulière, de sorte que, par 2.8 et la définition 2.5, la formule 2.9, appliquée aux diviseurs $s^*$ et $D'$ de $S_1$, s'écrit

\begin{equation}\tag{2.11.3}
(s^*,D) \;=\; \sum_{t' \text{ proche de } s} \mu_{t'}(D)\ [k(t'):k] \qquad .
\end{equation}
L'assertion (i) pour $s = t$ résulte de (2.11.2) et (2.11.3), et le cas général s'obtient en appliquant ce résultat au localisé en $t$ de $S(t)$ (2.3).

Désignons par $m$ l'idéal maximal qui définit $\{s\}$, et soit $n \geq 0$. On sait (SGA VII 3.5) que

\begin{equation}\tag{2.11.4}
p_*\ \underline{O}_{S_1}(-n\ s^*) = m^n \quad \text{et} \quad R^i p_*\ \underline{O}_{S_1}(-n\ s^*) = 0 \quad (i > 0) \qquad .
\end{equation}

%%%%%%%%%%  scan idx 27  |  folio 20  %%%%%%%%%%

On déduit de là, et de la suite exacte

\begin{equation*}
0 \longrightarrow \underline{O}_{S_1}(-n\ s^*) \longrightarrow \underline{O}_{S_1} \longrightarrow \underline{O}_{n\,s^*} \longrightarrow 0
\end{equation*}
que

\begin{equation}\tag{2.11.5}
p_*\ \underline{O}_{n\,s^*} \;=\; O_S/m^n \quad \text{et} \quad R^i p_*\ \underline{O}_{n\,s^*} \;=\; 0 \quad (i > 0)
\end{equation}
et que

\begin{equation}\tag{2.11.6}
\chi(\underline{O}_{n\,s^*}) \;=\; \chi(\underline{O}_S/m^n) \;=\; \frac{n(n+1)}{2} \ [k(s):k] \ .
\end{equation}
Faisons $n = 0$ dans (2.11.4). Quel que soit le complexe $K$ de faisceaux cohérents sur $S$, on a par la ``formule de projection''

\begin{equation*}
K \xrightarrow{\ \sim\ } Rp_*\ Lp^*K \qquad (\simeq K \overset{\mathbb{L}}{\otimes} Rp_*\ \underline{O}_{S_1}) \qquad .
\end{equation*}
En particulier, puisque $Lf^*\underline{O}_D = f^*\ \underline{O}_D = \underline{O}_{D^*}$, on a

\begin{equation}\tag{2.11.7}
p_*\ \underline{O}_{D^*} = \underline{O}_D \quad \text{et} \quad R^i p_*\underline{O}_{D^*} = 0 \quad (i>0) \qquad .
\end{equation}
Posons $E = \mu_s(D)\ s^*$ . D'après (2.11.1), (2.9.1) (i) b et (2.6.1) on a

\begin{equation}\tag{2.11.8}
\chi(\underline{O}_{D'\cap E}) = (D',E) = (D^* - E,E) = -(E,E) = -\mu_s(D)^2(s^*,s^*) = \mu_s(D)^2\ [k(s) : k] \qquad .
\end{equation}
La suite exacte

\begin{equation*}
0 \longrightarrow O_{D^*} \longrightarrow O_{D'} \oplus \underline{O}_E \longrightarrow \underline{O}_{D'\cap E} \longrightarrow 0
\end{equation*}
fournit par image directe, grâce à (2.11.7) des suites exactes

\begin{equation*}
0 \longrightarrow O_D \longrightarrow f_*\underline{O}_{D'} \oplus f_*\underline{O}_E \longrightarrow f_*\underline{O}_{D'\cap E} \longrightarrow 0
\end{equation*}
et

\begin{equation*}
0 \longrightarrow f_*\underline{O}_{D'}/\underline{O}_D \longrightarrow f_*\underline{O}_E \longrightarrow f_*\ \underline{O}_{D'\cap E} \longrightarrow 0 \qquad .
\end{equation*}

%%%%%%%%%%  scan idx 28  |  folio 21  %%%%%%%%%%

D'après (2.11.6) et (2.11.8), on a donc

\begin{equation}\tag{2.11.9}
\begin{aligned}
\chi(f_*\underline{O}_{D'}/\underline{O}_D) &= \frac{1}{2}\,\mu_s(D)(\mu_s(D)+1)[k(s): k] - \mu_s(D)^2\,[k(s) : k] \\
&= \frac{1}{2}\,\mu_s(D)(\mu_s(D)-1)[k(s) : k] \quad .
\end{aligned}
\end{equation}
Les schémas $D$ et $D'$ sont réduits, ne diffèrent qu'en $s$, et $D'$ est fini sur $D$. Ces schémas ont donc même normalisé, et l'invariant $\delta$ de la singularité de $D$ en $s$ est somme de $\chi(f_*\underline{O}_{D'}/\underline{O}_D)$ et de la somme des invariants $S$ des singularités de $D'$ en les points de $D' \cap s^*$. Prouvons 2.11 (ii) par récurrence sur $\delta$. Si $\mu_s(D) > 1$, l'hypothèse de récurrence s'applique à $D'$, et 2.11 (ii) résulte de (2.11.9). Si $\mu_s(D) = 1$, alors $(D',s^*) = 1$ et $D' \cap s^*$ est un diviseur régulier de $D'$ ; par (EGA IV 19.1.1), $D'$ est régulier ; par (2.11.9), $D = D'$ est lui-même régulier et 2.11 (ii) résulte de 2.8.

2.12. Soit $t \in S^*$ et $S'$ un transformé birationnel de $S$ qui domine $S_+(t)$ (2.3). On appellera \underline{transformé total} de $t$ sur $S'$, et on désignera encore par $t^*$, le transformé total sur $S'$ du diviseur $t^*$ de $S_+(t)$.

Soit $\nu$ une fonction à support fini de $S^*$ dans $\mathbf{N}$ et supposons que $\nu$ vérifie, quel que soit $t \in S^*$

\begin{equation}\tag{2.12.1}
\nu(t) \;\geq\; \sum_{t' \ \mathrm{proche\ de}\ t} \nu(t')\,[k(t') : k(t)]
\end{equation}
(inégalité des points proches). Soient $f : S' \longrightarrow S$ un transformé birationnel de $S$ qui domine les $S_+(t)$ pour $t$ dans le support de $\nu$, et $\nu^*$ le diviseur $\sum_{t \in S^*} \nu(t)\ t^*$ sur $S'$.


%%%%%%%%%%  scan idx 29  |  folio 22  %%%%%%%%%%

On désignera par $I(\nu)$ l'idéal de $\underline{O}_S$ image directe par $f$ de l'idéal de $\underline{O}_{S'}$, qui définit le diviseur $\nu^*$ :

\begin{equation*}
I(\nu) \;=\; f_*\underline{O}(-\nu^*) \quad .
\end{equation*}
D'après (2.11.4), cet idéal ne dépend pas du choix de $S'$. Les idéaux de $S$ obtenus par ce procédé sont ceux que Zariski [11] appelle ``complets''.

\underline{Théorème} 2.13. \underline{Sous les hypothèses précédentes}, \underline{on a}

\begin{equation*}
\chi(\underline{O}_S/I(\nu)) \;=\; \sum_{t \in S^*} \frac{\nu(t)(\nu(t)+1)}{2}\,[k(t) : k]
\end{equation*}
\underline{et}

\begin{equation*}
R^i f_*\,\underline{O}_{S'}(-\nu^*) \;=\; 0 \quad \underline{\mathrm{pour}}\ i > 0 \quad .
\end{equation*}
En vieux français, on exprime ceci, brièvement, en disant que : ``exiger qu'une courbe plane passe $n$ fois par un point $t$, projectif ou infiniment voisin, impose $\frac{n(n+1)}{2}$ conditions, en général indépendantes, aux coefficients de son équation''.

Prouvons tout d'abord le lemme

\underline{Lemme} 2.14. \underline{Soient} $p : S_1 \longrightarrow S$ \underline{l'éclaté de} $S$ \underline{le long de} $\{s\}$, $T$ \underline{un sous-schéma fermé de} $S_1$ , \underline{d'idéal} $q$, $D$ \underline{un diviseur}, $n$ \underline{un entier}, $J = O(-s^*)$ \underline{l'idéal de} $S_1$ \underline{qui définit le diviseur} $s^*$, $Y = \mathrm{Spec}(\underline{O}_{S_1}/J^n q)$, $D$ \underline{un diviseur régulier sur} $S$ \underline{et} $D'$ \underline{son transformé pur sur} $S_1$. \underline{On suppose} $T \cap s^*$ \underline{et} $T \cap D'$ \underline{finis}.

(i) \underline{Si} $\mathrm{lg}_{k(s)}\,(T \cap s^*) \leq n +1$, \underline{on a}

\begin{equation*}
R^i p_*\,(q.J^n) \;=\; 0 \quad \underline{\mathrm{pour}}\ i > 0 \ \underline{\mathrm{et}},\ \underline{\mathrm{pour}}\ T\ \underline{\mathrm{artinien}},
\end{equation*}
\begin{equation*}
\mathrm{lg}_{k(s)}(\underline{O}_S/p_*(q\,J^n)) \;=\; \frac{n(n+1)}{2} \;+\; \mathrm{lg}_{k(s)}\,(T)
\end{equation*}

%%%%%%%%%%  scan idx 30  |  folio 23  %%%%%%%%%%

(ii)  \underline{Si} $\mathrm{lg}_{k(s)} (T \cap s^*) \leq n$, \underline{on a}

\begin{equation*}
\mathrm{lg}_{k(s)} (D \cap \mathrm{Spec}(\underline{O}_S/p_*(q\,J^n))) \; = \; n + \mathrm{lg}_{k(s)} (D' \cap T) \quad .
\end{equation*}
Soit la suite exacte sur $S_1$

\begin{equation}\tag{2.14.1}
0 \longrightarrow J^n. q \longrightarrow J^n \stackrel{u}{\longrightarrow} J^n \otimes \underline{O}_T \longrightarrow 0 \quad .
\end{equation}
Pour prouver que $R^ip_* J^n.q = 0$  $(i > 0)$, il suffit par (2.11.4), de montrer que la flèche $p_*(u)$ est surjective. D'après le théorème des fonctions holomorphes, il suffit de prouver que les flèches

\begin{equation*}
u_k \; : \; p_*(J^n \otimes \underline{O}/J^k) \longrightarrow p_*(J^n \otimes \underline{O}_T \otimes \underline{O}/J^k)
\end{equation*}
sont surjectives.

On désignera par $\mathrm{Gr}_J$ un gradué associé pour la filtration $J$-adique. Le module $J^m/J^{m+1}$ est isomorphe à $\underline{O}_{s^*}(m)$ ; on a donc $R^1p_*J^m/J^{m+1} = 0$ et la flèche de $p_*(J^n/J^{n+k+1})$ dans $p_*(J^n/J^{n+k})$ est surjective. Pour prouver, par récurrence sur $k$, que les flèches $u_k$ sont surjectives, il suffit donc de prouver, la surjectivité des $v_k$ :

\begin{equation*}
v_k \; : \; p_*(J^n \otimes \mathrm{Gr}_J^k(\underline{O})) \longrightarrow p_*(J^n \otimes \mathrm{Gr}_J^k(\underline{O}_T)) \quad (k \geq 0) \quad .
\end{equation*}
Le module $\mathrm{Gr}_J^k(\underline{O}_T)$ est un quotient de $J^k/J^{k+1} \otimes \underline{O}_T$ ; il est donc de longueur sur $k(s)$ au plus $n+1$. Le noyau $N$ de la flèche surjective

\begin{equation*}
v_k' \; : \; J^n \otimes \mathrm{Gr}_J^k(\underline{O}) \sim \underline{O}_{s^*}(\,n+k) \longrightarrow J^n \otimes \mathrm{Gr}_J^k(\underline{O}_T)
\end{equation*}
est donc un faisceau inversible sur la droite projective $s^*$, de degré $\geq -1$. On a $R^1p_* N = 0$ et $v_k = p_*(v_k')$ est surjectif.

Si $T$ est artinien, de la suite exacte (2.14.1), on déduit maintenant que

\begin{equation*}
\mathrm{lg}_{k(s)} (p_*J^n/p_*(J^nq)) = \mathrm{lg}_{k(s)} (\underline{O}_T) \quad , 
\end{equation*}

%%%%%%%%%%  scan idx 31  |  folio 24  %%%%%%%%%%

 et 2.14 (i) résulte de (2.11.4) et (2.11.6).

Soit $I_{D'\cap T}$ l'idéal (monogène de $\underline{O}_{D'}$ qui définit $D'\cap T$. Sous les hypothèses (ii), la suite exacte

\begin{equation*}
0 \longrightarrow I_{D'\cap T} \longrightarrow \underline{O}_{D'\cap T} \longrightarrow 0_T \longrightarrow 0
\end{equation*}
montre que $D'\cap T$ vérifie l'hypothèse (i).

La projection $p$ établit un isomorphisme entre $D$ et $D'$ ; par cet isomorphisme, à l'idéal de $\underline{O}_D$ qui définit le sous-schéma $D \cap \mathrm{Spec}(O_S/p_*(q\,J^n))$ de $D$ correspond l'idéal de $D'$ engendré par les sections globales de $q\,J^n$ sur $S_1$. Cet idéal est contenu dans l'idéal de $D'$ qui définit $D'\cap Y$. Il lui est en fait égal, car toute section de $\underline{O}_{D'\cup Y}$ se prolonge en une section de $\underline{O}_{S_1}$ : si $I_{D'\cup Y}$ est l'idéal qui définit $D'\cup Y$, l'obstruction ou prolongement se trouve dans $R^1p_*(I_{D'\cup Y})$. On a

\begin{equation*}
\begin{array}{rl} I_{D'\cup Y} & = \underline{O}(-D')\cap q\,J^n = (\underline{O}(-D')\cap J^n) \cap q\,J^n = \underline{O}(-D')J^n \cap q\,J^n \\[2mm] & = (\underline{O}(-D')\cap q). J^n \end{array}
\end{equation*}
et, d'après (i) appliqué à $D'\cup T$, il n'y a pas d'obstruction.

On a donc

\begin{equation}\tag{2.14.2}
\mathrm{lg}_{k(s)} (D\cap \mathrm{Spec}(\underline{O}_S/p_*(q\,J^n))) = \mathrm{lg}_{k(s)} (D'\cap Y) \quad .
\end{equation}
De l'isomorphie

\begin{equation*}
J^n \otimes \underline{O}_{D'}/q.\underline{O}_{D'} \stackrel{\sim}{\longrightarrow} \underline{O}_{D'}.J^n/\underline{O}_{D'}.q.J^n
\end{equation*}
on déduit que la suite suivante est exacte

\begin{equation*}
0 \longrightarrow J^n \otimes_{D'\cap T} \longrightarrow \underline{O}_{D'\cap Y} \longrightarrow \underline{O}_{D'\cap\ n\ s^*} \longrightarrow 0 \quad ,
\end{equation*}
de sorte que 


%%%%%%%%%%  scan idx 32  |  folio 25  %%%%%%%%%%

\begin{equation}\tag{2.14.3}
\mathrm{lg}_{k(s)}\ (D' \cap Y) = n\ (D',s^*) + \mathrm{lg}_{k(s)}(D' \cap T)
\end{equation}
Puisque $(D',s^*) = 1$, 2.14 (ii) résulte de (2.14.2) et (2.14.3).

Prouvons maintenant, par récurrence sur le nombre d'éléments dans le support de $\nu$, les assertions de 2.13 et l'assertion suivante :

Si $D$ est un diviseur régulier sur $S$, alors,

\begin{equation*}
\chi(D \cap \mathrm{Spec}(\underline{O}_S/I(\nu)) = \sum_{t\ \mathrm{sur}\ D} \nu(t)\ [k(t) : k]
\end{equation*}
Si $\nu = 0$, alors $\nu^* = 0$, $I(\nu) = \underline{O}_S$ et 2.13 et (2.14.4) sont triviaux (la seconde assertion de (2.13) résulte d'une application itérée de (2.11.4) pour $n = 0$). Sinon, on a $\nu(s) > 0$, et il existe un $S$-morphisme $g : S' \longrightarrow S_1$. Pour chaque point fermé $a$ de $S_1$, soient $S_{1a}$ le localisé de $S_1$ en $a$, $\nu_a$ la restriction de $\nu$ à $S_{1a}^*$ , $S'_a$ le localisé $S'_a = S' \times_{S_1} S_{1a}$, $I(\nu_a) = g_* \underline{O}_{S'}(-\nu_a^*)$, $T_a = \mathrm{Spec}(\underline{O}_{S_1}/I(\nu_a))$, $T = \cup\ T_a$ (réunion finie disjointe) et $q$ l'idéal de $T$. On a, avec les notations de 2.14

\begin{equation}\tag{2.14.5}
g_* \underline{O}(-\nu^*) = q.J^{\nu(s)}
\end{equation}
L'hypothèse de récurrence s'applique aux $(S_{1a}, S'_a, \nu_a)$. On a donc

\begin{equation}\tag{2.14.6}
R^i g_* \underline{O}_{S'}(-\nu^*) = 0 \ \mathrm{pour}\ i > 0
\end{equation}
\begin{equation}\tag{2.14.7}
\chi(T) = \sum_{t \in S^* \setminus \{s\}} \frac{\nu(t)\ (\nu(t)+1)}{2} [k(t) : k]
\end{equation}
et, d'après (2.14.4) appliqué au diviseur $s^*$ de $S_1$ et l'hypothèse (2.12.1)

\begin{equation*}
\mathrm{lg}_{k(s)}\ (T \cap s^*) = \sum_{t\ \mathrm{proche\ de}\ s} \nu(t)\ [k(t) : k(s)] \leq \nu(s)
\end{equation*}

%%%%%%%%%%  scan idx 33  |  folio 26  %%%%%%%%%%

Cette inégalité nous permet d'appliquer 2.14 ; de (2.14.5), on tire

\begin{equation*}
R^i p_* \ g_* \ \underline{O}(-\nu^*) = 0 \ \mathrm{pour}\ i > 0
\end{equation*}
\begin{equation*}
\mathrm{lg}_{k(s)}(\underline{O}_S/p_* \ g_* \ \underline{O}(-\nu^*)) = \frac{\nu(s)(\nu(s)+1)}{2} + \mathrm{lg}_{k(s)}\ (T)
\end{equation*}
de sorte que les assertions (2.13) résultent de (2.14.6) et (2.14.7). Enfin, 2.14 (ii) fournit

\begin{equation*}
\mathrm{lg}_{k(s)}\ (D \cap \mathrm{Spec}\ (\underline{O}_S/I(\nu)) = \nu(s) + \mathrm{lg}_{k(s)}\ (D' \cap T) \ ,
\end{equation*}
soit par (2.14.4) appliqué à $D'$,

\begin{equation*}
\begin{aligned}
\chi(D \cap \mathrm{Spec}(\underline{O}_S/I(\nu)) &= \nu(s)\ [k(s) : k] + \sum_{t\ \mathrm{sur}\ D} \nu(t)\ [k(t) : k] \\
&= \sum_{t\ \mathrm{sur}\ D} \nu(t)\ [k(t) : k]
\end{aligned}
\end{equation*}
Ceci achève la démonstration de 2.13.

Des bribes de ce qui précèdent se généralisent en dimension supérieure. Soit $S$ le spectre d'un anneau local régulier de dimension $n$ de point fermé $s$, $k(s)$ étant une extension finie de $k$. On peut encore considérer l'ensemble $I$ des schémas $f : S' \longrightarrow S$ déduit de $S$ par éclatements successifs de points fermés et poser comme en 2.1

\begin{equation*}
S^* = \varinjlim_{s' \in I} \ \mathrm{Irr}(f^{-1}(s)
\end{equation*}
$S^*$ s'appellera encore l'ensemble des \underline{points de} $S$ \underline{infiniment voisins de} $s$.

Si $X$ est un sous-schéma fermé de $S$, et si $t \in S^*$, on définit comme en 2.8 la multiplicité de $X$ en $t$.

Soient $f : S' \longrightarrow S$ un composé de morphismes d'éclatement de points fermés, $D$ un diviseur sur $S$, $D^* = f^* D$ son transformé total sur $S'$,
%%%%%%%%%%  scan idx 34  |  folio 27  %%%%%%%%%%
$E$ une courbe sur $S$, $E'$ son transformé pur sur $S'$ et supposons que $D \cap E$ soit réduit à $s$. La formule de projection donne encore

\begin{equation}\tag{2.16}
(D,E) \;=\; (D^*,E')
\end{equation}
Procédant comme en 2.9, on prouve à l'aide de (2.16) que

\begin{equation}\tag{2.17}
(D,E) \;=\; \sum_{t \in S^*} \mu_t(D) \, \mu_t(E) \, [k(t) : k] \qquad .
\end{equation}
Définissons les points proches comme en 2.5. Procédant comme en 2.11, on vérifie que

\begin{equation}\tag{2.18}
\mu_t(E) \;=\; \sum_{t' \text{ proche de } t} \mu_{t'}(E) \, [k(t') : k] \qquad .
\end{equation}
Par contre, la formule 2.11 (ii) devient fausse en général (voir les exemples 2.25).

2.19. Soit $S$ une surface projective et régulière sur un corps $k$. On désigne par $S^*$ la somme disjointe, étendue aux points fermés de $S$

\begin{equation*}
S^* \;=\; \coprod_{s \in S} \operatorname{Spec} \, (\underline{O}_{s,S})^* \qquad ,
\end{equation*}
et on appelle $S^*$ l'ensemble des \underline{points infiniment voisins des} points de $S$.

Soit $I$ l'ensemble des surfaces déduites de $S$ par éclatements successifs de points fermés. L'ensemble $I$ est cofinal dans l'ensemble des modèles réguliers du corps des fonctions de $S$, ordonné par domination. Pour $i \in I$, on désigne par $S_i$ la surface correspondantes, munie de $p_i : S_i \to S$. On désignera par $\mathrm{Div}^{\pm}(S_i)$ le groupe des diviseurs, positifs ou non, sur $S_i$.

\underline{Définition} 2.20. (i) \underline{Le groupe des} ``courbes'' \underline{sur} $S$ \underline{est l'ensemble limite inductive}

\begin{equation*}
\varinjlim_{i \in I} \mathrm{Div}^{\pm}(S_i)
\end{equation*}
(ii) \underline{La catégorie des} ``faisceaux inversibles'' \underline{sur} $S$ \underline{est la limite inductive}
%%%%%%%%%%  scan idx 35  |  folio 28  %%%%%%%%%%
(SGA 4 VI 6) \underline{des catégories de faisceaux inversibles sur les} $(S_i)_{i \in I}$ .

Ces objets sont des invariants birationnels de $S$.

\underline{Exemple} 2.21. (i) Chaque ``courbe'' $C$ sur $S$ définit un ``faisceau inversible'' $\underline{O}(C)$ sur $S$.

(ii) Chaque diviseur $C$ sur $S$ définit une ``courbe'' sur $S$, qu'on désignera par $C^*$.

(iii) Si $t \in S^*$, on définit comme en 2.3 la surface $S_+(t)$ et le diviseur $t^*$ sur $S_+(t)$ ; on désigne encore par $t^*$ la ``courbe'' définie par $t^*$.

(2.22) On vérifie facilement qu'une ``courbe'' $C$ (resp. un ``faisceau inversible'' $\mathcal{L}$) s'écrit d'une et d'une seule façon sous la forme

\begin{equation*}
C \;=\; C_o^* \;-\; \sum_{t \in S^*} i_t \, t^* \qquad \text{resp.} \qquad \mathcal{L} \;=\; \mathcal{L}_o \otimes \underline{O} \left( -\sum_{t \in S^*} i_t \, t^* \right)
\end{equation*}
où $C_o$ (resp. $\mathcal{L}_o$) est un diviseur (resp. un faisceau inversible) sur $S$ et où les $i_t$ sont des nombres entiers presque tous nuls. On dit que $C$ (resp. $\mathcal{L}$) est la courbe $C_o$, (resp. le faisceau inversible $\mathcal{L}_o$) affectée de multiplicités virtuelles $i_t$.

Le complété $\mathrm{Div}^{\pm} (S) \times \mathbb{Z}^{S^*}$ du groupe $\mathrm{Div}^{\pm} (S) \times \mathbb{Z}^{(S^*)}$ des ``courbes'' sur $S$ est encore un invariant birationnel de $S$ ; on l'appelle le groupe des ``courbes généralisées'' sur $S$. Sur $S$, une ``courbe généralisée'' $C$ se représente par un diviseur $C_o$, affecté d'un nombre pouvant être infini de multiplicités virtuelles.

On définit de même les ``faisceaux inversibles généralisés'' sur $S$.


%%%%%%%%%%  scan idx 36  |  folio 29  %%%%%%%%%%

\underline{Exemple} 2.23. Supposons $k$ parfait. Alors, le faisceau inversible $\Omega^2_{S/k}$ , affecté d'une multiplicité virtuelle $-1$ en chaque point $t \in S^*$, et un invariant birationnel de $S$, soit $\Omega^2_{S/k}{}^*$ .

\underline{Exemple} 2.24.  Soit $f : S_1 \longrightarrow S_2$ un morphisme dominant de surfaces projectives régulières sur $k$. On définit sans difficulté l'image inverse d'une "courbe" ou d'un "faisceau inversible" sur $S_2$, en passant à un modèle de $S_2$ où cette "courbe" ou "faisceau inversible" est diviseur ou faisceau inversible usuel. Par passage à la limite (je n'ai pas vérifié si ce passage à la limite était possible), on définit de même l'image réciproque d'une "courbe généralisée", ou d'un"faisceau inversible généralisé" . Supposons $k$ parfait et $f$ génériquement étale. La \underline{différente} $\mathfrak{d}(f)$ de $f$ est alors le diviseur positif sur $S_1$ qui vérifie

\begin{equation*}
f^* \Omega^2_{S_2/k} (\mathfrak{d}(f)) \overset{\sim}{\longrightarrow} \Omega^2_{S_1/k} \qquad .
\end{equation*}
Cette différente n'est pas un invariant birationnel de $f$ ; pour obtenir un invariant birationnel, il faut considérer la "courbe généralisée" de coïncidences

\begin{equation*}
\mathfrak{d}^*(f) = \mathfrak{d}(f) + \sum_{t \in S_1^*} t^* - f^* ( \sum_{t \in S_2^*} t^*)
\end{equation*}
soit, brièvement

\begin{equation*}
\mathfrak{d}^*(f) = \Omega^2_{S_1/k}{}^* - f^* \Omega^2_{S_2/k}{}^* \qquad .
\end{equation*}
Cette "courbe généralisée" à un support contenu dans le fermé de $S_2$ où $f$ n'est pas étale.


%%%%%%%%%%  scan idx 37  |  folio 30  %%%%%%%%%%

(2.25) \underline{Exemple}.

Soient $k$ un corps algébriquement clos et $E_n$ l'espace affine à $n$ dimension sur $k$. Soit $C$ une courbe réduite dans $E_n$, présentant une singularité à l'origine. A cette singularité on associe son "\underline{diagramme d'Enriques}" (cf. [3]),qui est le graphe de sommets les points infiniment voisins de $O$ \underline{sur} $C$ (2.4 et 2.15), chaque sommet étant joint par un trait à ceux qui lui sont proches (2.5 et 2.15). Ce diagramme détermine les multiplicités de ces points sur $C$ (2.11 et 2.18) et, pour $n = 2$, l'invariant $\delta$ (2.10 et 2.11) de la singularité de $C$ à l'origine.

Pour les singularités planes les plus simples, on trouve :

\noindent Dessin \hspace{5em} équation \hspace{5em} $\delta$ \hspace{5em} diagramme (avec multiplicités)

\noindent [dessin : un trait horizontal droit] \hspace{4em} $y = 0$ \hspace{4em} $0$

\begin{center}
\begin{tikzcd}[column sep=large]
\overset{1}{\bullet} \arrow[r,dash] & \overset{1}{\bullet} \arrow[r,dash] & \overset{1}{\bullet} \arrow[r,dash] & \overset{1}{\bullet} \arrow[r,dash] & \overset{1}{\bullet} \arrow[r,dash] & \overset{1}{\bullet} \arrow[r,dash] & \overset{1}{\bullet} & \dots
\end{tikzcd}
\end{center}
% STRUCTURE: TABLEAU, LIGNE 1 (équation $y = 0$, $\delta = 0$). Colonne « Dessin » : un unique trait
%   horizontal droit tracé à main levée (aucun texte). Colonne « diagramme (avec multiplicités) » :
%   chaîne horizontale de 7 sommets marqués par un point sur un trait continu. NŒUDS (7) : N1, N2, N3,
%   N4, N5, N6, N7, alignés horizontalement, chacun surmonté du chiffre 1 dactylographié (label
%   AU-DESSUS du nœud). ARÊTES (6), toutes non orientées, trait simple sans pointe, horizontales, sans
%   label : N1—N2, N2—N3, N3—N4, N4—N5, N5—N6, N6—N7. À droite de N7, cinq points dactylographiés sur la
%   ligne de base (marque de continuation « ..... ») : ce n'est PAS un sommet ; rendu par la cellule
%   finale \dots. Aucune flèche, aucun label d'arête, aucun trait pointillé, aucune double flèche.
\noindent [dessin : deux traits droits qui se croisent, en forme de $\times$] \hspace{4em} $xy = 0$ \hspace{4em} $1$

\begin{center}
\begin{tikzcd}[column sep=large, row sep=small]
{} & \overset{1}{\bullet} \arrow[r,dash] & \overset{1}{\bullet} \arrow[r,dash] & \overset{1}{\bullet} \arrow[r,dash] & \overset{1}{\bullet} \arrow[r,dash] & \overset{1}{\bullet} & \dots \\
\overset{2}{\bullet} \arrow[ur,dash] \arrow[dr,dash] & {} & {} & {} & {} & {} & {} \\
{} & \overset{1}{\bullet} \arrow[r,dash] & \overset{1}{\bullet} \arrow[r,dash] & \overset{1}{\bullet} \arrow[r,dash] & \overset{1}{\bullet} \arrow[r,dash] & \overset{1}{\bullet} & \dots
\end{tikzcd}
\end{center}
% STRUCTURE: TABLEAU, LIGNE 2 (équation $xy = 0$, $\delta = 1$). Colonne « Dessin » : deux traits droits
%   sécants formant une croix (un X), tracés à main levée (aucun texte). Colonne « diagramme » : NŒUDS
%   (11) — R (racine, à gauche, à mi-hauteur), surmonté du chiffre 2 ; branche SUPÉRIEURE H1, H2, H3,
%   H4, H5 alignés horizontalement, chacun surmonté du chiffre 1 (labels AU-DESSUS de la ligne) ;
%   branche INFÉRIEURE B1, B2, B3, B4, B5 alignés horizontalement, chacun surmonté du chiffre 1 (labels
%   également AU-DESSUS de la ligne inférieure). ARÊTES (10), toutes non orientées, trait simple sans
%   pointe, sans label : R—H1 (oblique montante vers la droite), R—B1 (oblique descendante vers la
%   droite), H1—H2, H2—H3, H3—H4, H4—H5 (horizontales), B1—B2, B2—B3, B3—B4, B4—B5 (horizontales). Après
%   H5 : cinq points dactylographiés (continuation) ; après B5 : cinq points dactylographiés
%   (continuation). Ces points ne sont pas des sommets. Aucune flèche orientée, aucun trait pointillé,
%   aucun label d'arête.
\noindent [dessin : un trait horizontal et une parabole tangente à ce trait] \hspace{4em} $y (y - x^2) = 0$ \hspace{4em} $2$

\begin{center}
\begin{tikzcd}[column sep=large, row sep=small]
{} & {} & \overset{1}{\bullet} \arrow[r,dash] & \overset{1}{\bullet} \arrow[r,dash] & \overset{1}{\bullet} \arrow[r,dash] & \bullet & {} \\
\overset{2}{\bullet} \arrow[r,dash] & \overset{2}{\bullet} \arrow[ur,dash] \arrow[dr,dash] & {} & {} & {} & {} & {} \\
{} & {} & \underset{1}{\bullet} \arrow[r,dash] & \underset{1}{\bullet} \arrow[r,dash] & \underset{1}{\bullet} \arrow[r,dash] & \bullet & \dots
\end{tikzcd}
\end{center}
% STRUCTURE: TABLEAU, LIGNE 3 (équation $y(y - x^2) = 0$, $\delta = 2$). Colonne « Dessin » : un trait
%   horizontal droit et une parabole (ouverte vers le haut) tangente à ce trait en un point, tracés à
%   main levée (aucun texte). Colonne « diagramme » : NŒUDS (10) — A (extrême gauche) surmonté du
%   chiffre 2 ; B (à sa droite, même hauteur) surmonté du chiffre 2, c'est le point de branchement ;
%   branche SUPÉRIEURE H1, H2, H3, H4 alignés horizontalement, H1, H2, H3 surmontés chacun du chiffre 1
%   (labels AU-DESSUS), H4 (extrémité droite du trait, marquée par un point plus épais) SANS label ;
%   branche INFÉRIEURE B1, B2, B3, B4 alignés horizontalement, B1, B2, B3 portant chacun le chiffre 1 EN
%   DESSOUS de la ligne, B4 (extrémité droite, point plus épais) SANS label. ARÊTES (9), toutes non
%   orientées, trait simple sans pointe, sans label : A—B (horizontale), B—H1 (oblique montante), B—B1
%   (oblique descendante), H1—H2, H2—H3, H3—H4, B1—B2, B2—B3, B3—B4. Après H4 : AUCUN point de
%   continuation. Après B4 : DEUX points dactylographiés seulement (« .. »), rendus par la cellule
%   finale \dots. Aucune flèche orientée, aucun trait pointillé.
\noindent [dessin : une courbe à point de rebroussement, la pointe tournée vers la gauche] \hspace{4em} $y^2 = x^3$ \hspace{4em} $1$

\begin{center}
\begin{tikzcd}[column sep=large, row sep=small]
\overset{2}{\bullet} \arrow[r,dash] \arrow[drr,dash] & \overset{1}{\bullet} \arrow[dr,dash] & {} & {} & {} \\
{} & {} & \overset{1}{\bullet} \arrow[r,dash] & \overset{1}{\bullet} \arrow[r,dash] & \overset{1}{\bullet} & \dots
\end{tikzcd}
\end{center}
% STRUCTURE: TABLEAU, LIGNE 4 (équation $y^2 = x^3$, $\delta = 1$). Colonne « Dessin » : une courbe
%   présentant un point de rebroussement (cusp), la pointe tournée vers la gauche, les deux branches
%   s'ouvrant vers la droite, tracée à main levée (aucun texte). Colonne « diagramme » : NŒUDS (5) — A
%   (extrême gauche, niveau supérieur) surmonté du chiffre 2 ; B (à droite de A, même niveau supérieur)
%   surmonté du chiffre 1 ; C (niveau inférieur, sous et à droite de B) surmonté du chiffre 1 ; D
%   (niveau inférieur, à droite de C) surmonté du chiffre 1 ; E (niveau inférieur, à droite de D)
%   surmonté du chiffre 1. ARÊTES (5), toutes non orientées, trait simple sans pointe, sans label : A—B
%   (horizontale courte), A—C (longue oblique descendante, du niveau supérieur au niveau inférieur), B—C
%   (oblique descendante plus courte) — ces trois arêtes forment un TRIANGLE A, B, C ; puis C—D et D—E
%   (horizontales, niveau inférieur). Après E : cinq points dactylographiés (continuation), rendus par
%   la cellule finale \dots ; ce n'est pas un sommet. Aucune flèche orientée, aucun trait pointillé.
Si $n > 2$, l'invariant $\delta$ n'est plus déterminé par ce diagramme. Ainsi, la réunion de 3 droites concourantes coplanaires (resp. non coplanaires) a un invariant $\delta = 3$ (resp. $\delta = 2$), alors que le diagramme d'Enriques est le même. J'ignore si ce phénomène se présente aussi pour des courbes unibranches.

Voici deux exemples de diagrammes, relatifs au cas $n = 3$, qui ne peuvent se présenter pour des courbes planes (cf. 2.5).

On donne successivement les équations de la courbe, les équations
%%%%%%%%%%  scan idx 38  |  folio 31  %%%%%%%%%%
de la transformée pure après éclatement de l'origine, des équations paramétriques de la courbe, l'invariant $\delta$ et le diagramme d'Enriques.

\begin{equation}\tag{2.25.1}
\left\{
\begin{array}{l}
\left\{\begin{array}{l} y^{2} = xz^{4} \\ z^{3} = x^{2} \end{array}\right. \;\;,\;\;
\left\{\begin{array}{l} (y/z)^{2} = (x/z)^{3} \\ z \;\;\;\;\;\; = (x/z)^{2} \end{array}\right. \;\;,\;\;
\left\{\begin{array}{l} x = u^{6} \\ y = u^{7} \\ z = u^{4} \end{array}\right. \;\;,\;\; \delta = 5 \;\;,
\end{array}
\right.
\end{equation}
\begin{equation}\tag{2.25.1}
\begin{tikzcd}[column sep=small, row sep=small]
4 \arrow[r, dash] \arrow[rrd, dash] \arrow[rrdd, dash] & 2 \arrow[rd, dash] \arrow[rdd, dash] & {} & {} & {} & {} \\
{} & {} & 1 \arrow[d, dash] & {} & {} & {} \\
{} & {} & 1 \arrow[rd, dash] & {} & {} & {} \\
{} & {} & {} & 1 \arrow[rd, dash] & {} & {} \\
{} & {} & {} & {} & 1 & {} \\
{} & {} & {} & {} & {} & \cdots
\end{tikzcd}
\end{equation}
% STRUCTURE: Enriques diagram printed to the right of the tag (2.25.1), inside the same large outer
%   brace as the equations. 6 dot-vertices, each carrying a printed multiplicity label, plus a trailing
%   ellipsis of three dots that is NOT joined by any line. All edges are plain undirected line segments
%   (no arrowheads, no labels, no hooks, no dashes). Vertices: v1 label 4 (top left, label printed
%   above-left of the dot); v2 label 2 (top, to the right of v1, label above); v3 label 1 (below and to
%   the right of v2, label to the right of the dot); v4 label 1 (almost directly below v3, label to the
%   right); v5 label 1 (below and right of v4, label above-right); v6 label 1 (below and right of v5,
%   label above-right). Edges, 8 in total, all dash/undirected: v1--v2 (horizontal, rightwards); v1--v3
%   (down-right, shallow); v1--v4 (down-right, steep); v2--v3 (down-right); v2--v4 (down-right, steeper
%   than v2--v3); v3--v4 (vertical, downwards); v4--v5 (down-right, drawn slightly curved); v5--v6
%   (down-right, drawn slightly curved). The segments v1--v3 and v2--v4 cross each other (verified by
%   pixel scan: the two lines merge into a single dark run at the crossing column). After v6 the drawing
%   stops; three separate dots trail off down-right.
\begin{equation}\tag{2.25.2}
\left\{
\begin{array}{l}
\left\{\begin{array}{l} y^{2}z = x^{3} \\ z^{3} \;\; = y^{2} \end{array}\right. \;\;,\;\;
\left\{\begin{array}{l} (y/z)^{2} = (x/z)^{3} \\ z \;\;\;\;\;\; = (y/z)^{2} \end{array}\right. \;\;,\;\;
\left\{\begin{array}{l} x = u^{8} \\ y = u^{9} \\ z = u^{6} \end{array}\right. \;\;,\;\; \delta = 10 \;\;,
\end{array}
\right.
\end{equation}
\begin{equation}\tag{2.25.2}
\begin{tikzcd}[column sep=small, row sep=small]
6 \arrow[r, dash] \arrow[rrd, dash] \arrow[rrdd, dash] \arrow[rrddd, dash] \arrow[rrdddd, dash] & 2 \arrow[rd, dash] \arrow[rdd, dash] & {} & {} & {} & {} \\
{} & {} & 1 \arrow[d, dash] & {} & {} & {} \\
{} & {} & 1 \arrow[d, dash] & {} & {} & {} \\
{} & {} & 1 \arrow[d, dash] & {} & {} & {} \\
{} & {} & 1 \arrow[rd, dash] & {} & {} & {} \\
{} & {} & {} & 1 \arrow[rd, dash] & {} & {} \\
{} & {} & {} & {} & 1 & {} \\
{} & {} & {} & {} & {} & \cdots
\end{tikzcd}
\end{equation}
% STRUCTURE: Enriques diagram printed to the right of the tag (2.25.2), inside the same large outer
%   brace as the equations. 8 dot-vertices, each carrying a printed multiplicity label, plus a trailing
%   ellipsis of three dots that is NOT joined by any line. All edges are plain undirected line segments
%   (no arrowheads, no labels, no hooks, no dashes). Vertices: w1 label 6 (top left, label above); w2
%   label 2 (top, to the right of w1, label above); w3 label 1 (lower right of w2, label to its right);
%   w4 label 1 (directly below w3, label to its right); w5 label 1 (directly below w4, label to its
%   right); w6 label 1 (directly below w5, label to its right); w7 label 1 (below and right of w6, label
%   above-right); w8 label 1 (below and right of w7, label above-right). Edges, 12 in total, all
%   dash/undirected: w1--w2 (horizontal, rightwards); w1--w3, w1--w4, w1--w5, w1--w6 (a fan of four
%   segments going down-right from w1 with increasing slope — verified by pixel scan: 5 distinct dark
%   runs in the columns between w1 and w2, and 6 runs in the columns between w2 and the vertical chain);
%   w2--w3 (down-right); w2--w4 (down-right, steeper); w3--w4, w4--w5, w5--w6 (the single near-vertical
%   line running downwards through w3, w4, w5, w6); w6--w7 (down-right); w7--w8 (down-right). After w8
%   the drawing stops; three separate dots trail off down-right.
2.26. Le théorème 2.27 ci-dessous était énoncé sous la forme d'une conjecture fausse dans la première version (diffusée par l'IHES) de cet exposé. Il doit beaucoup à des conversations avec D.S. Rim, dont il englobe certains résultats (D.S. Rim, torsion differentials and deformations, theorem 2.7 (i)).


%%%%%%%%%%  scan idx 39  |  folio 32  %%%%%%%%%%

Soit  $C$  une courbe projective réduite sur un corps algébriquement clos  $k$ , et  $s \in C(k)$. On désignera par

$A_{s}$ : l'anneau local de  $C$  en  $s$  (ou son complété, cela reviendrait au même pour ce qui suit) ;

$\widetilde{A}_{s}$ : le normalisé de  $A_{s}$ ;

$\delta(s):\dim_{k}(\widetilde{A}_{s}/A_{s})$  ;

$m_{o}(s)$ : la dimension du schéma versel formel  $T$  des déformations de $\mathrm{Spec}(A_{s})$.

Pour  $E$  une composante irréductible de  $T$ , on posera encore $m_{o}^{E}(s) = \dim(E)$.

Les modules  $\widetilde{\tau}$  des dérivations de  $\widetilde{A}_{s}$  et  $\tau$  des dérivations de  $A_{s}$  sont sans torsion ; ils s'injectent dans le module des dérivations du corps des fractions de  $A_{s}$ , pour  $A$  intègre, et dans le module des dérivations de l'anneau total des fractions de  $A_{s}$  (un produit de corps) dans le cas général. De plus, via ces inclusions,  $\tau$  et  $\widetilde{\tau}$  sont commensurables : $\tau/\tau\cap\widetilde{\tau}$  et  $\widetilde{\tau}/\tau\cap\widetilde{\tau}$  sont de dimension finie sur  $k$ .

On pose

\begin{equation*}
m_{1}(s) = [\widetilde{\tau} : \tau] \mathop{=}\limits_{\mathrm{dfn}} \dim_{k}(\widetilde{\tau}/\tau \cap \widetilde{\tau}) - \dim_{k}(\tau/\tau\cap\widetilde{\tau}) \quad .
\end{equation*}
\underline{Théorème} 2.27.  \underline{Soit}  $E$  \underline{une composante irréductible du schéma formel versel}  $T$  \underline{des déformations de}  $A_{s}$ . \underline{Si}  "\underline{la déformation de}  $A_{s}$  \underline{au-dessus du point générique}  $e$  \underline{de}  $E$  \underline{est lisse}", \underline{i.e.} \underline{si l'image du sous-schéma formel de non lissité} (\underline{défini par un idéal jacobien} : VI 5), \underline{fini sur}  $T$ , \underline{ne contient pas}  $e$ , \underline{alors}

\begin{equation*}
3\ \delta(s) = m_{o}^{E}(s) + m_{1}(s) \quad .
\end{equation*}

%%%%%%%%%%  scan idx 40  |  folio 33  %%%%%%%%%%

La validité de 2.27 ne dépend que du complété de $A_s$ . Remplaçant $(C,s)$ par $(C',s)$ formellement isomorphe à $(C,s)$ en $s$ , on peut supposer que les conditions suivantes sont vérifiées

(a)\quad $s$ est le seul point singulier de $C$ ;

(b)\quad $C$ n'a pas d'automorphismes infinitésimaux.

La condition (b) nous garantit que la foncteur des déformations de $C$ est proreprésentable. Il est même effectivement pro-représentable, car il n'y a pas d'obstruction à relever un faisceau inversible (ample) donné sur $C$ . Soit $f : X_o \longrightarrow T_o$ la déformation universelle de $C$ et $t_o$ le point fermé de $T_o$ : $C \stackrel{\sim}{\relbar\joinrel\relbar} X_{ot_o}$.

\underline{Lemme} 2.28. \quad \underline{Il existe un schéma de type fini sur} $k$ $T_1$, $t_1 \in T_1(k)$ \underline{et} $f : X_1 \longrightarrow T_1$, \underline{propre et plat,} \underline{de fibre} $C$ \underline{en} $t_1$ \underline{et tel que pour tout} $u \in T_1(k)$, \underline{le complété de} $T_1$ \underline{en} $u$ \underline{soit un schéma formel de modules pour la fibre} $X_{1u}$ .

Voici une démonstration directe de ce corollaire de résultats généraux d'Artin. Soit $\mathcal{L}$ un faisceau inversible très ample sur $C$ , tel que $H^1(C,\mathcal{L}) = 0$, et $\mathbb{P} = \mathbb{P}(H^o(C,\mathcal{L}))$. On a $v : C \hookrightarrow \mathbb{P}$ . Soit $T_2$ le schéma de Hilbert, classifiant les sous-variétés de $\mathbb{P}$ de même polynôme de Hilbert que $v(C) \subset \mathbb{P}$ . Soit $t_2 \in T_2$ correspondant à $v(C)$. Le complété $\hat{T}_2$ de $T_2$ en $t_2$ est lisse sur $S_o$ , le fibré tangent vertical de $\hat{T}_2/T_o$ se déduisant du sous-fibré $L$ du fibré tangent de $T_2$ image de $\mathrm{Lie}(\underline{\mathrm{Aut}}(\mathbb{P}))$ pour l'action de $\underline{\mathrm{Aut}}(\mathbb{P})$ sur $T_2$. Au voisinage de $t_2$, soit $f_1 \ldots f_n$ $n = \dim \underline{\mathrm{Aut}}(\mathbb{P})$ équations définissant un sous-schéma $T'_1$ de $T_2$ passant par $t_2$ et de complété en $t_2$ étale sur $T_o$ . On prend pour $T_1$ un voisinage de $t$ dans $T'_1$ où le fibré tangent soit transverse à $L$ .


%%%%%%%%%%  scan idx 41  |  folio 34  %%%%%%%%%%

2.29. Au morphisme de foncteur

\begin{equation*}
(\text{déformations de } C) \longrightarrow (\text{déformations de } A_s)
\end{equation*}
correspond un morphisme $p : T_o \longrightarrow T$ . D'après VI 2.10 \qquad , $p$ est lisse. Soit $t$ le point fermé de $T$ . L'espace tangent de Zariski de $p^{-1}(t)$ en $t_o$ classifie les déformations localement triviales de $C$ sur $k[\varepsilon]$ ($\varepsilon^2 = 0$). Posant $E_o = p^{-1}(E)$, on a donc

\begin{equation}\tag{2.29.1}
\dim(E_o) = \dim(E) + \dim H^1(C, T^1_{C/k}) \qquad .
\end{equation}
Puisque $X_o/T_o$ est plat, les fibres de $X_o/T_o$ ont toutes le même genre arithmétique $p$ ($1-p = \chi(\mathcal{O})$). Par hypothèse, la fibre au point générique de $E_o$ est lisse ; le schéma formel des modules d'une courbe propre et lisse étant lisse de dimension $3p-3$, on déduit de 2.28 que

\begin{equation}\tag{2.29.2}
\dim(E) \;=\; 3p-3 \qquad .
\end{equation}
Soit $\widetilde{C}$ la normalisée de $C$, de genre arithmétique $\widetilde{p} = p - \delta(s)$. On a

\begin{equation}\tag{2.29.3}
\begin{aligned}\dim\, H^1(C,T^1_{C/k}) \;&=\; -\,\chi(C,T^1_{C/k}) \;=\; \\ &=\; -\,\chi(\widetilde{C},T^1_{\widetilde{C}/k}) + m_1(s) \;=\; -[3(p-\delta(s))-3] + m_1(s) \qquad .\end{aligned}
\end{equation}
Les formules (2.29.1) (2.29.2) (2.29.3) prouvent 2.27.

\underline{Conjecture} 2.30. \quad \underline{Toute courbe propre réduite a une déformation lisse.}

Cette conjecture implique, par 2.27 et 2.28, que le schéma formel versel des déformations de $A_s$ est toujours équidimensionnel, de dimension $3\delta(s) - m_1(s)$.


%%%%%%%%%%  scan idx 42  |  folio 35  %%%%%%%%%%

2.31. Soit $C$ une courbe propre réduite sur $k$, de genre arithmétique $p$ . On désignera par

$M_o$ : la dimension du schéma formel versel $T$ des déformations de $C$ ;

$M_1$ : la dimension de l'algèbre de Lie de $\underline{\mathrm{Aut}}(C)$.

Pour $E$ une composante irréductible de $T$, on notera $M_o^E$ la dimension de $E$ .

L'analogue global de 2.27 est le suivant.

\underline{Proposition} 2.32. \underline{Si la déformation de} $C$ \underline{au-dessus du point générique de} $E$ \underline{est lisse, on a}

\begin{equation*}
3p-3 = M_o^E - M_1 \quad .
\end{equation*}
Soit $X \subset C(k)$ un ensemble fini de points lisses de $C$ , tel que tout champ de vecteur sur $C$, s'annulant en tout point $x \in X$, soit nul. Le foncteur des déformations de $(C,X)$ est proreprésentable, et vérifie une variante de 2.28 . Le schéma formel des modules de $(C,X)$ est lisse sur $T$ , de dimension relative $\# X - M_1$ , et de l'analogue de 2.28 on déduit que

\begin{equation*}
M_o + ( \# X - M_1) \;=\; 3p - 3 + \# X \; ,
\end{equation*}
d'où 2.32.

Supposons $k$ de \underline{caractéristique} $0$ . On a alors $M_1 = \dim \underline{\mathrm{Aut}}(C)$. Voici une interprétation analogue de $m_1(s)$.


%%%%%%%%%%  scan idx 43  |  folio 36  %%%%%%%%%%

\underline{Lemme} 2.33. \underline{Si} $k$ \underline{est de caractéristique $0$, toute dérivation de} $A_s$ \underline{se prolonge en une dérivation de} $A_s^{\sim}$ .

Si $D$ est une dérivation de $A_s$ , $\exp(Dt) = \sum D^n t^n/n!$ est un automorphisme du complété en $(s,0)$ de $A_s \otimes_k k[t]$. Par transport de structure, cet automorphisme agit sur le normalisé $(A_s^{\sim} \otimes k[t])^{\wedge}$ de cet anneau. Réduisant mod $t^2$ , on trouve un automorphisme $1 + Dt$ de $A_s^{\sim\wedge} \otimes_k k[t]/(t^2)$, et $D$ se trouve ainsi prolongé en une dérivation de $A_s^{\sim\wedge}$ . Puisqu'il se prolonge aussi au corps des fractions de $A_s$ (anneau total des fractions si $A_s$ n'est pas intègre), il se prolonge à $A_s^{\sim}$ .

Supposons $s$ singulier. Un argument analogue montre alors que le champ de vecteur sur $\mathrm{Spec}(A_s^{\sim})$ défini par $D$ s'annule en tout point au-dessus de $s$ .

Soit $c$ un idéal de $A_s^{\sim}$ contenu dans le conducteur de $A_s^{\sim}/A_s$ . Soit $G_c^{\sim} = \underline{\mathrm{Aut}}(A_s^{\sim}/c)$ et soit $G_c$ le sous-groupe de $G_c^{\sim}$ qui stabilise la sous-algèbre $A_s/c$ . Pour $d \subset c$ , $G_d^{\sim}$ stabilise $c/d$ et s'envoie sur $G_c^{\sim}$ (remarquer que $g \in G_c^{\sim}$ est déterminé par $g(t)$, pour $t$ une uniformisante, et que $g(t)$ est n'importe quelle uniformisante). L'espace homogène pointé $G_c^{\sim}/G_c$ est donc indépendant de $c$ . On note $G^{\sim}/G$.

Soit $\tau_o^{\sim}$ le $A_s^{\sim}$-module des champs de vecteurs sur $\mathrm{Spec}(A_s^{\sim})$ qui s'annulent en les points au-dessus de $s$ . Pour $s$ singulier, on

\begin{equation*}
\dim_k(\tau^{\sim}/\tau_o^{\sim}) \;=\; \mbox{nombre de branches} \;;
\end{equation*}
\begin{equation*}
\tau_o^{\sim}/\tau \;\stackrel{\sim}{\longrightarrow}\; \mbox{espace tangent à l'origine de } G^{\sim}/G \;.
\end{equation*}

%%%%%%%%%%  scan idx 44  |  folio 37  %%%%%%%%%%

\underline{Proposition} 2.34. \underline{En caractéristique} $0$ \underline{et pour} $s$ \underline{singulier}, \underline{l'entier} $m_1(s)$ \underline{est somme du nombre de branches et de la dimension de} $\widetilde{G}/G$ .


\underline{BIBLIOGRAPHIE}

[1] \quad P. Defrise. Etude locale des correspondances rationnelles entre surfaces algébriques - Thèse d'agrégation - Bruxelles 1949.

[2] \quad P. Deligne et D. Mumford. The irreducibility of the space of curves of a given genus - Publ. Math. IHES \underline{36} (1969) p. 75-110.

[3] \quad Enriques - Chisini. Lezioni sulla teoria geometrica delle equazioni et delle funzioni algebriche - ed Zanichelli - Bologne 1924.

[4] \quad R. Hartshorne. Residues and duality - Lectures Notes in Math. \underline{20} - Springer 1966.

[5] \quad S. Lichtenbaum. Curves over discrete valuation rings - Am. J. of Math. 15 n° 2 April 1968. p. 380-405.

[6] \quad I. Manin. Rational surfaces over perfect fields - Publ. Math. IHES n° 30. p. 55-114.

[7] \quad D. Mumford. The topology of normal singularities of an algebraic surfaces and a criterion for simplicity - Publ. Math. IHES n° 9 p. 5-22.

[8] \quad A. Néron. Modèles minimaux des variétés abéliennes sur les corps locaux et globaux - Publ. Math. n° 21.


%%%%%%%%%%  scan idx 45  |  folio 38  %%%%%%%%%%

[9] \quad M. Raynaud. Spécialisation du foncteur de Picard. Critères numériques de représentabilité - CR Acad. Sci. Paris t. 264 p. 1001 - 1004 - Juin 1967.

[10] \quad Saint-Donat. Sur un théorème de G. Castelnuovo et F. Enriques - Thèse de 3° cycle - Lyon.

[11] \quad O. Zariski. On polynomial ideals defined by infinitely near base points Am. J. Math. vol. 60-1938 p. 151.

[12] \quad D.G. Northcott. On the genus formula for plane curves. J. London Math. Soc. 30 1955 - p. 376-382.

[13] \quad D.G. Northcott. Abstract dilatations and infinitely near points. Proc. Cambridge Philos. Soc. 52 1956 - p. 178-197.

La notion de ``point infiniment voisin'' est discutée dans l'appendice par Lipman au chapitre II de

[14] \quad O. Zariski. Algebraic surfaces - Ergebnisse \underline{61}, Springer Verlag.
